\documentclass[11pt]{article}
\usepackage[margin=1in]{geometry}
\usepackage[T1]{fontenc}
\usepackage{lmodern}
\usepackage{microtype}
\usepackage{amsmath,amssymb,amsthm,mathtools}
\usepackage{booktabs,longtable,tabularx,array}
\usepackage{enumitem}
\usepackage{xcolor}
\usepackage{listings}
\usepackage{fancyhdr}
\usepackage{hyperref}
\usepackage[nameinlink,noabbrev]{cleveref}
\hypersetup{
  colorlinks=true,
  linkcolor=blue!45!black,
  citecolor=green!35!black,
  urlcolor=blue!55!black,
  pdftitle={The Barnes G-function in intuitionistic primitive recursive arithmetic},
  pdfauthor={Research report prepared for Vladimir Reshetnikov},
  pdfsubject={Constructive quantitative proofs and Lean formalization plan}
}
\pagestyle{fancy}
\fancyhf{}
\lhead{Barnes $G$ in iPRA}
\rhead{16 November 2026}
\cfoot{\thepage}
\setlength{\headheight}{14pt}
\setlength{\parindent}{0pt}
\setlength{\parskip}{0.55em}
\setlist{nosep,leftmargin=2em}
\allowdisplaybreaks

\newtheorem{theorem}{Theorem}[section]
\newtheorem{lemma}[theorem]{Lemma}
\newtheorem{proposition}[theorem]{Proposition}
\newtheorem{corollary}[theorem]{Corollary}
\newtheorem{definition}[theorem]{Definition}
\newtheorem{remark}[theorem]{Remark}
\newtheorem{contract}[theorem]{Finite contract}
\newtheorem{audit}[theorem]{Source audit}
\newtheorem{roadmap}[theorem]{Formalization step}

\newcommand{\N}{\mathbb N}
\newcommand{\Z}{\mathbb Z}
\newcommand{\Q}{\mathbb Q}
\newcommand{\R}{\mathbb R}
\newcommand{\C}{\mathbb C}
\newcommand{\dd}{\,\mathrm d}
\newcommand{\ii}{\mathrm i}
\newcommand{\e}{\mathrm e}
\newcommand{\Log}{\operatorname{Log}}
\newcommand{\Cl}{\operatorname{Cl}_2}
\newcommand{\sfac}{\operatorname{sf}}
\newcommand{\hyp}{\operatorname{H}}
\newcommand{\BallEq}{\operatorname{BallEq}}
\newcommand{\code}[1]{\texttt{#1}}
\newcommand{\statusA}{\textbf{A}}
\newcommand{\statusB}{\textbf{B}}
\newcommand{\statusC}{\textbf{C}}
\newcommand{\statusD}{\textbf{D}}

\lstset{
  basicstyle=\ttfamily\small,
  columns=fullflexible,
  breaklines=true,
  frame=single,
  backgroundcolor=\color{black!2}
}

\title{\textbf{The Barnes $G$-function in intuitionistic primitive recursive arithmetic}\\
\large Quantitative constructive proofs, source audit, and a Lean-to-iPRA formalization plan}
\author{Research report prepared for Vladimir Reshetnikov}
\date{16 November 2026}

\begin{document}
\maketitle

\begin{abstract}
We reconstruct, in a form suitable for intuitionistic primitive recursive
arithmetic (iPRA), essentially every mathematical identity on the Wikipedia and
MathWorld pages for the Barnes $G$-function.  The central device is not a
classical completed real or complex universe, but a family of finite statements:
for every fixed primitive-recursive name, compact rational domain, requested
precision and branch-separation certificate, a bounded checker verifies a
rational enclosure.  The Weierstrass product is given an explicit uniform tail
bound; this yields the entire function, its zero multiplicities, the functional
equation, integer and finite-product identities, Taylor and derivative series,
integral and reflection formulas, multiplication and duplication formulas,
quarter- and half-integer values, Hurwitz-zeta and hyperfactorial relations,
and order two of infinite type.  Euler--Maclaurin and an explicit Barnes-$G$
remainder integral give a fixed-order, fixed-sector asymptotic schema.

Two source errors are corrected.  The asymptotic characterization uses
$\Gamma(n)^x$, not $\Gamma(x)^n$, and the pure-imaginary absolute-value formula
has the prefactor $|y|$, not $y$.  We also separate three levels of evidence:
the human mathematical proof developed here, executable finite regression
checks, and the still-to-be-completed compilation into ProveIt's object theory
$[]\vdash_{\mathrm{PRAAx}}\varphi$.  The accompanying Lean files specify all
necessary definitions, certificate interfaces and target formulas, but are not
presented as a completed formal proof.
\end{abstract}

\tableofcontents
\newpage

\section{Statement of the problem and exact scope}

The two target pages mix four very different kinds of assertion.

\begin{enumerate}
\item There are exact finite identities, such as
\(G(n)=0!1!\cdots(n-2)!\) and the finite product of gamma values.
\item There are convergent infinite constructions, such as the canonical
product and the Taylor series.
\item There are analytic continuation statements involving a logarithm,
square root, quotient or path integral.  These are not meaningful until a
zero-free domain and a branch have been fixed.
\item There are asymptotic and uniqueness statements that quantify over a
precision, a sector, a differentiability order, or a coded function.
\end{enumerate}

A proof ``in iPRA'' must therefore say what finite arithmetical sentence replaces
each assertion.  It is not enough to copy a classical proof whose objects are
arbitrary Cauchy sequences or arbitrary functions.  Conversely, iPRA is not
limited to tiny finite computations.  A fixed approximation algorithm, an
explicit primitive-recursive modulus, and quantifier-free induction are enough
to express all of the finite estimates used below.

\subsection{Four status levels}\label{sec:scope-status}

Every entry in the coverage ledger is assigned one of the following levels.

\begin{description}
\item[\statusA: exact finite theorem.]  The claim reduces to natural-number,
integer or rational identities and inequalities.  The report gives a direct
primitive-recursive induction.  These are the first targets to compile into
\code{Prf PRAAx}.

\item[\statusB: quantitative analytic theorem.]  The report gives a finite
algorithm and an explicit error bound on a certified compact domain.  For every
fixed code and precision, the resulting statement is first-order arithmetic.
The remaining work is implementation of the checkers and their object-language
derivations.

\item[\statusC: branch- or representation-dependent theorem.]  The classical
formula is valid only after a continuation path, a branch of logarithm or
square root, or a normalization of a negative-order polygamma has been chosen.
The report supplies the missing data and proves the corresponding finite
contract.

\item[\statusD: corrected source statement.]  The displayed source formula is
false or misleading as written.  We identify the defect, prove the corrected
formula, and retain the original in the audit only.
\end{description}

None of these letters means ``already kernel-checked in Lean.''  A completed
object-iPRA theorem has the more specific form
\[
  []\;\vdash[\mathrm{PRAAx}]\;\varphi,
\]
where the terms in \(\varphi\) are built from the primitive-recursive symbols of
ProveIt's \code{ConciseIPRA.PR}.  The present package supplies the mathematics,
the exact target formulas and the dependency order, but does not claim that all
analytic targets have already been constructed as such derivations.

\subsection{Pinned repository foundation}

At the pinned ProveIt commit \cite{ProveIt}, iPRA is presented as first-order
intuitionistic logic with equality, a symbol for every primitive-recursive
function, its defining equations, successor axioms, and induction for
quantifier-free formulas.  The checked library has three relevant layers:

\begin{enumerate}
\item \code{ConciseIPRA} defines \code{PR}, terms, formulas, the axiom family
\code{PRAAx}, derivations \code{Prf}, and the bridge \code{pra\_iff\_native} to
the small native scalar calculus.
\item \code{IPRADerivations} supplies checked arithmetic: semiring laws,
order, bounded sums and products, division with remainder, pairing,
simultaneous recursion and polynomial normalization.
\item the \code{gamma\_ipra} packet gives a broad quantitative plan for Gamma,
including finite Euler products, explicit constants, integration, reflection,
multiplication and Stirling estimates.  Its own status statement correctly says
that the full analytic Lean development and the compilation into the object
theory remain work.
\end{enumerate}

This report reuses the third layer as a source of contracts, not as an assumed
theorem oracle.  Whenever a Barnes proof invokes Gamma, the formal dependency
is a named finite Gamma contract: an evaluator, a modulus, an apartness result,
or a fixed-sector Stirling remainder.  Such a contract can later be discharged
by the Gamma packet's intended modules.

\subsection{The theorem collection obtained}

The main conclusions are the following.

\begin{theorem}[Coverage theorem]\label{thm:coverage}
For every fixed primitive-recursive input name and every requested precision,
the following Barnes-$G$ statements admit an iPRA formulation by bounded
certificates and quantifier-free induction.

\begin{enumerate}[label=(\roman*)]
\item the canonical product, with a uniform compact-disc tail;
\item existence as an entire function, real conjugation, and the zero at
\(0,-1,-2,\ldots\) with multiplicities \(1,2,3,\ldots\);
\item normalization and the functional equation;
\item all positive-integer values, the superfactorial identity, the
hyperfactorial identity, and finite Gamma products;
\item the Taylor series for \(\log G(1+z)\), with explicit tails for every
closed subdisc of \(|z|<1\);
\item the logarithmic derivative, derivative, log-Gamma integral, and the
improper-integral representation on \(\Re z>-1\);
\item Kinkelin reflection, its Clausen forms, and the half-shift tangent form;
\item the squared absolute-value products, including the corrected
pure-imaginary formula;
\item the general multiplication formula, duplication, quarter values,
half-integer values, and the integer coefficient formula for all odd
half-integers;
\item the Hurwitz-zeta derivative representation, the negative-order
polygamma representation after normalization, and the analytic K-function
relation;
\item the full fixed-order Barnes asymptotic expansion on every fixed sector
separated from the negative real axis, with an explicit remainder bound;
\item order two and infinite type, with explicit integer-radius estimates;
\item the Vign\'eras characterization and the asymptotic characterization,
read as schemas for fixed coded functions with effective smoothness or
relative-error moduli.
\end{enumerate}
\end{theorem}

The proof occupies the rest of the report.  Its constructive content is not a
new axiom of real analysis.  It is the repeated pattern
\[
  \text{finite approximant}
  \quad+\quad
  \text{primitive-recursive error budget}
  \quad+\quad
  \text{bounded verifier}.
\]

\subsection{Two errors caught before formalization}

\begin{audit}[Asymptotic characterization]\label{audit:char}
The Wikipedia page presently displays
\(f(x+n)\sim\Gamma(x)^n n^{\binom{x}{2}}f(n)\).  Setting \(x=1\) contradicts
the functional equation, because the left side is exactly
\(\Gamma(n)f(n)\), while the displayed right side is \(f(n)\).  The cited
higher-order Bohr--Mollerup theorem has
\[
 f(x+n)\sim \Gamma(n)^x n^{\binom{x}{2}}f(n).
\]
We prove this corrected formula in \cref{sec:characterization}.
\end{audit}

\begin{audit}[Pure-imaginary modulus]\label{audit:modulus}
The Wikipedia page gives a factor \(y\) on the right of a formula for
\(|G(\ii y)|\) and says that \(y\) is arbitrary real.  The right side would be
negative for \(y<0\).  Conjugating the product gives a factor \(y^2\) in the
squared formula, and hence the unsquared factor is \(|y|\).  The corrected
identity is proved in \cref{sec:modulus}.
\end{audit}

These are not cosmetic matters.  The finite checker rejects both incorrect
statements immediately: the first at the exact input \(x=1\), the second at any
negative rational \(y\).  This illustrates why source auditing and formal
normalization should precede large proof-assistant developments.

\section{The constructive language: names, balls, and finite queries}\label{sec:ipra}

The classical notation \(z\in\C\) is convenient in the proof, but it cannot be
the primitive ontology of iPRA.  This section gives a uniform translation from
that notation to natural-number data.

\subsection{Exact finite arithmetic}

Integers are represented by pairs \((a,b)\in\N^2\), denoting \(a-b\).
Rationals are represented by triples \((a,b,d)\) with \(d>0\), denoting
\((a-b)/d\).  Equality and order are checked by cross multiplication; gcd
normalization is optional.  A complex rational is a pair of rationals.  All
operations used below---addition, multiplication, conjugation, powers, finite
sums and products---are primitive recursive.

For analytic error estimates it is useful to avoid square roots.  We use the
one-norm
\[
 \|(x,y)\|_1=|x|+|y|
\]
on complex rationals.  It majorizes the Euclidean norm and is submultiplicative
up to the exact inequality
\(\|zw\|_1\le \|z\|_1\|w\|_1\).  Thus every ball calculation can be reduced to
rational arithmetic.  When a theorem is conventionally stated with the
Euclidean modulus, the comparison
\[
 |z|\le \|z\|_1\le 2\max(|\Re z|,|\Im z|)
\]
changes only a harmless fixed guard in the precision budget.

Factorials, superfactorials, hyperfactorials, Bernoulli numbers of fixed index,
finite harmonic sums, and every finite product appearing in this report are
primitive recursive.  Division is never used as an unguarded operation in an
object theorem.  A quotient statement \(a/b=c\) is first represented by
\(a=bc\), plus a certificate that \(b\ne0\) if the quotient itself is needed.

\subsection{Fixed fast names rather than a real-number sort}

\begin{definition}[Fast name]
A \emph{complex fast name} is a fixed primitive-recursive function
\(a:\N\to\Q^2\) satisfying
\[
 \|a(p+1)-a(p)\|_1\le 2^{-p}.
\]
Its intended value is the unique complex number in all nested tail balls.
Two names \(a,b\) are equal when, for every \(p\),
\[
 \|a(p+2)-b(p+2)\|_1\le 2^{-p}.
\]
\end{definition}

This definition is metatheoretic shorthand.  iPRA does not quantify over all
functions \(a\).  For each concrete primitive-recursive symbol, the regularity
inequality is a theorem; for each \(p\), equality at precision \(p\) is a
finite rational inequality.  The statement ``the two represented numbers are
equal'' is the schema consisting of all those precision theorems.

\begin{definition}[Apartness]
A name \(a\) is quantitatively apart from zero if finite data \((m,c)\), with
\(c>0\), certify
\[
  \|a(m)\|_1-2^{-m+1}\ge c2^{-m}.
\]
The certificate supplies a primitive-recursive lower bound on the true
modulus and hence a safe precision budget for reciprocal, logarithm and
quotient computations.
\end{definition}

Unlike the classical disjunction \(z=0\lor z\ne0\), apartness is computational:
it tells the evaluator how far it is from a singularity.  All quotient formulas
in later sections are first proved in cross-multiplied form and converted only
under an apartness hypothesis.

\subsection{Compact domains and branch paths}

A compact input domain is specified by rational data, for example
\[
 D(c,R)=\{z:\|z-c\|_1\le R\}.
\]
A branch of \(\Log\) is specified by a rational polygonal path from a positive
base point to the input, together with a positive rational lower bound on the
distance of every segment from the chosen cut and from zero.  The continuation
algorithm subdivides the path until each logarithmic increment lies in a disk
where the atanh series
\[
 \Log w=2\sum_{j=0}^{M-1}\frac{u^{2j+1}}{2j+1}+\text{tail},
 \qquad u=\frac{w-1}{w+1},
\]
has a certified geometric tail.  The same pattern defines rational powers,
square roots and continuously chosen arguments.

Thus a formula such as
\(\Log G(1-z)-\Log G(1+z)=\cdots\) is not asserted on an unspecified global
branch.  Its finite input contains paths for both sides or, more economically,
a single zero-free simply connected domain with a branch tree.  Exponentiating
then produces a branch-independent identity.

\subsection{Finite sums, products, integrals and derivatives}

A continuous named function on a rational rectangle is represented by a fixed
primitive-recursive evaluator
\(E(q,p)\) returning a dyadic ball, plus a modulus of uniform continuity.
A Riemann integral at precision \(p\) is computed on a mesh selected from that
modulus.  The verifier checks:

\begin{enumerate}
\item every sample ball;
\item the mesh bound;
\item the oscillation bound on each cell;
\item any endpoint or infinite-tail estimates;
\item the final rational interval containment.
\end{enumerate}

Derivatives are represented either by a named derivative with a Taylor
remainder certificate, or by symmetric difference quotients with an explicit
modulus.  Holomorphy on a compact set is most efficient here as a Taylor-ball
contract: around each point of a finite rational cover there are coefficients
and a geometric remainder.  The canonical product supplies those coefficients
by uniformly convergent derivative series.

Improper integrals are always split at finite rational cutoffs.  For example,
\(
\int_0^\infty f(t)\dd t
\)
is replaced by
\(
\int_0^{2^{-M}}+\int_{2^{-M}}^T+\int_T^\infty
\).
The first and last pieces receive analytic majorants; only the middle piece is
quadrature.  All cutoffs are primitive-recursive functions of the requested
precision and domain parameters.

\subsection{Finite-query normal form}

The following organization matches the reusable framework in the repository.

\begin{definition}[FQNF contract]
A theorem family is in \emph{finite-query normal form} when it is described by
primitive-recursive predicates and functions
\[
 \operatorname{Pre}(q),\quad B(q),\quad
 \operatorname{Verify}(c,q),\quad \operatorname{Goal}(c,q),
\]
such that:
\begin{enumerate}
\item from \(\operatorname{Pre}(q)\), a primitive-recursive producer returns a
certificate \(c\le B(q)\);
\item \(\operatorname{Verify}(c,q)\) is quantifier-free;
\item iPRA proves
\(\operatorname{Verify}(c,q)\to\operatorname{Goal}(c,q)\);
\item every loop invariant in the producer is quantifier-free, so its
correctness is proved by the induction rule of \code{PRAAx}.
\end{enumerate}
\end{definition}

For Barnes $G$, a typical query is
\[
 q=(\text{input code},\text{precision},\text{domain radius},
    \text{branch path},\text{apartness data}).
\]
The certificate records a product cutoff, Taylor lengths, Gamma and exponential
subcertificates, and a final dyadic ball.  A backward demand pass allocates
precision to the subcomputations.  The soundness proof is a finite telescope of
local ball inclusions.

\begin{contract}[Canonical-product query]\label{contract:canonical-query}
For a natural radius \(R\), precision \(p\), and an input name certified by
\(|z|\le R\), choose
\[
 N\ge 2R,
 \qquad
 \frac{2R^3}{3N}\le 2^{-p-3}.
\]
The certificate contains the first \(N\) factors, the named constants
\(\pi,\gamma\), and exponential/logarithm subcertificates.  The verifier checks
the displayed rational inequality and the finite factor balls.  The tail
lemma in \cref{lem:canonical-tail} then proves the requested enclosure.
\end{contract}

\subsection{Constructive logical discipline}

No proof below uses excluded middle for an arbitrary real predicate.  All case
splits are on decidable natural-number or rational comparisons.  Positivity is
used only when accompanied by a rational lower bound.  Limits are never invoked
without a modulus.  ``For sufficiently large \(n\)'' is replaced by an explicit
primitive-recursive threshold.  A fixed asymptotic order \(M\) is a numeral (or
an external schema parameter), so the finite list of Bernoulli coefficients is
primitive recursive without a universal evaluator for all primitive-recursive
functions.

This is exactly the strength needed for the arguments: they are constructive
analysis carried out by finite rational certificates, not classical analysis
smuggled into arithmetic by an undefined real-number predicate.

\section{The canonical product, convergence, zeros, and entirety}\label{sec:product}

Put
\begin{equation}\label{eq:Fdef}
 F(z):=G(1+z)
 =(2\pi)^{z/2}
 \exp\!\left(-\frac{z+(1+\gamma)z^2}{2}\right)
 \prod_{k=1}^{\infty}E_k(z),
\end{equation}
where
\begin{equation}\label{eq:Ek}
 E_k(z):=\left(1+\frac zk\right)^k
          \exp\!\left(-z+\frac{z^2}{2k}\right).
\end{equation}
This is the product displayed by DLMF, Wikipedia and MathWorld.  Its form is
already suggestive of a constructive proof: after the linear and quadratic
terms of \(k\Log(1+z/k)\) are cancelled, the tail is summable like \(k^{-2}\).

\subsection{The one estimate that drives the construction}

For \(|z|\le R\) and \(k\ge 2R\), use the principal Taylor branch of the
logarithm in \(|z/k|\le 1/2\) and set
\[
 q_k(z):=k\Log\left(1+\frac zk\right)-z+\frac{z^2}{2k}.
\]
The first two Taylor terms cancel exactly:
\begin{equation}\label{eq:qseries}
 q_k(z)=\sum_{m=3}^{\infty}
       \frac{(-1)^{m+1}z^m}{m k^{m-1}}.
\end{equation}

\begin{lemma}[Uniform canonical-factor tail]\label{lem:canonical-tail}
If \(|z|\le R\), \(N\ge 2R\), and \(N\ge1\), then
\begin{align}
 |q_k(z)|&\le \frac{2R^3}{3k^2}\qquad(k>N),\label{eq:qbound}\\
 \sum_{k>N}|q_k(z)|&\le \frac{2R^3}{3N}.
 \label{eq:qtail}
\end{align}
Consequently, for \(M>N\),
\begin{equation}\label{eq:product-cauchy}
 \left|\prod_{k=N+1}^{M}E_k(z)-1\right|
 \le \exp(B_N)B_N,
 \qquad B_N=\frac{2R^3}{3N}.
\end{equation}
\end{lemma}

\begin{proof}
From \cref{eq:qseries} and \(R/k\le1/2\),
\[
 |q_k(z)|
 \le \sum_{m=3}^{\infty}\frac{R^m}{m k^{m-1}}
 \le \frac{R^3}{3k^2}\sum_{j=0}^{\infty}\left(\frac Rk\right)^j
 \le \frac{2R^3}{3k^2}.
\]
The elementary integral comparison
\[
 \sum_{k=N+1}^{\infty}\frac1{k^2}
 \le \int_N^{\infty}\frac{dt}{t^2}=\frac1N
\]
is proved arithmetically by the telescoping inequality
\(
1/k^2\le1/(k-1)-1/k
\)
for \(k\ge2\), after clearing denominators.  This proves
\cref{eq:qtail}.  Finally
\(
\prod E_k=\exp(\sum q_k)
\)
and
\(
|\e^w-1|\le \e^{|w|}|w|
\), whose finite Taylor proof has an explicit remainder.  Applying it to the
finite sum and then passing through the Cauchy modulus gives
\cref{eq:product-cauchy}.
\end{proof}

Every inequality in this proof has a division-free version.  A convenient
primitive-recursive cutoff is
\begin{equation}\label{eq:product-cutoff}
 N(R,p):=\max\!\left(2R,\left\lceil\frac{2^{p+5}R^3}{3}\right\rceil\right).
\end{equation}
Then \(B_N\le2^{-p-5}\).  On this range \(B_N\le1\), so
\(\e^{B_N}B_N\le3B_N\le2^{-p-3}\), leaving guard bits for the
finite factors and the prefactor.  The ceiling is implemented by bounded
integer division, and the verifier checks the cross-multiplied inequality
rather than trusting the formula for \(N\).

\subsection{Derivative tails}

Uniform convergence of values alone gives a continuous function.  To obtain a
constructive holomorphic contract, differentiate the logarithmic tail.  Direct
calculation gives
\begin{align}
 q_k'(z)&=\frac1{1+z/k}-1+\frac zk
        =\frac{(z/k)^2}{1+z/k},\label{eq:qprime}\\
 q_k''(z)&=-\frac{1}{k(1+z/k)^2}+\frac1k
        =\frac{2z/k+z^2/k^2}{k(1+z/k)^2},\label{eq:qsecond}\\
 q_k^{(r)}(z)&=\frac{(-1)^{r-1}(r-1)!}
 {k^{r-1}(1+z/k)^r}\qquad(r\ge3).
 \label{eq:qhigher}
\end{align}
Thus, when \(|z|\le R\) and \(k\ge2R\),
\begin{align}
 |q_k'(z)|&\le \frac{2R^2}{k^2},\label{eq:qprime-bound}\\
 |q_k''(z)|&\le \frac{12R}{k^2},\label{eq:qsecond-bound}\\
 |q_k^{(r)}(z)|&\le
 2^r(r-1)!\,k^{1-r}\qquad(r\ge3).
 \label{eq:qhigher-bound}
\end{align}
For each fixed derivative order the tails have explicit moduli.  At orders
\(r\ge3\), the geometric/integral estimate is even faster.  These bounds give
a finite Taylor-ball certificate on every rational disk: evaluate finitely many
factors and sum finitely many derivative tails.  No appeal to a nonconstructive
``uniform limit theorem'' is necessary; that theorem is re-proved here in the
only quantitative form required.

\begin{theorem}[Constructive existence and entirety]\label{thm:entire}
The partial products in \cref{eq:Fdef} form, on every rational compact disk, a
fast uniformly Cauchy family together with every fixed derivative.  Their limit
\(F\) has a primitive-recursive evaluation modulus and compatible Taylor-ball
certificates.  Hence \(G(z)=F(z-1)\) is entire in the constructive sense of
\cref{sec:ipra}.
\end{theorem}

\begin{proof}
Fix an integer radius \(R\).  The first \(2R\) factors and the prefactor are
finite combinations of polynomials and exponentials, for which the Gamma packet
already specifies rational-ball evaluators.  Apply
\cref{lem:canonical-tail,eq:qprime-bound,eq:qsecond-bound,eq:qhigher-bound} to
the remaining factors.  Given precision \(p\) and derivative order \(r\), a
primitive-recursive maximum of the corresponding cutoffs makes every tail at
most \(2^{-p-c_r}\).  The finite product rule and exponential error lemmas
combine them into a Taylor ball.  Compatibility of overlapping balls follows
from the same value-tail estimate.  All loops are bounded by the cutoff, so the
certificate producer is primitive recursive; all local checks are rational
inequalities, so its soundness proof uses quantifier-free induction.
\end{proof}

\subsection{Normalization and real conjugation}

At \(z=0\), every factor in \cref{eq:Fdef} is one and the prefactor is one.
Therefore
\begin{equation}\label{eq:G1}
 G(1)=F(0)=1.
\end{equation}
This is exact, not merely a limiting statement.

Every coefficient and constant in the product is real.  Finite partial products
therefore satisfy
\(
F(\overline z)=\overline{F(z)}
\)
exactly.  The common tail modulus permits passage to the named limits:
\begin{equation}\label{eq:conj}
 G(\overline z)=\overline{G(z)}.
\end{equation}
In iPRA the proof at precision \(p\) simply evaluates both sides using the same
cutoff; the finite centers are conjugate and the two radius bounds coincide.

\subsection{Zeros and their multiplicities}

\begin{theorem}[Zero divisor]\label{thm:zeros}
The zeros of \(G\) are exactly
\[
 0,-1,-2,\ldots,
\]
and the zero at \(1-m\) has multiplicity \(m\) for every \(m\ge1\).
\end{theorem}

\begin{proof}
For \(F(z)=G(1+z)\), the finite factor
\((1+z/m)^m\) vanishes to order \(m\) at \(z=-m\).  Its exponential multiplier
never vanishes.  Every other finite factor is nonzero there, as is the
prefactor.  Choose \(N>2m+2\).  On a sufficiently small rational disk around
\(-m\), all factors with \(k>N\) satisfy \(|z/k|<1/2\); their logarithmic tail
converges to a finite value, so their product is quantitatively apart from zero.
The remaining finitely many nonvanishing factors have a rational lower bound on
a smaller disk.  Therefore \(F(z)=(z+m)^m U_m(z)\), where \(U_m\) is analytic
and apart from zero on that disk.

Conversely, at a point that is not a negative integer, choose a rational disk
avoiding the finitely many possible zeros among the head factors.  Their product
has a positive rational lower bound, and the tail exponential is nonzero.  Thus
there is no other zero.  Translating by one gives the stated divisor of \(G\).
\end{proof}

The factorization proof is fully finitary: for a fixed integer \(m\), polynomial
division extracts \((z+m)^m\) from the \(m\)-th factor, while all lower bounds
are supplied by bounded searches over rational boxes.  This proves the integer
zero statements on both target pages and strengthens them by multiplicity.

\subsection{Why the product is the right iPRA definition}

One could instead define \(G\) from its functional equation plus a uniqueness
principle.  That route postpones rather than removes the analytic work: the
uniqueness condition involves third derivatives or an asymptotic modulus.  The
canonical product is preferable as the foundational iPRA definition because
its convergence is controlled by the single rational inequality
\cref{eq:qtail}.  Functional equations, reflection and special values then
become theorems with independently checkable finite certificates.

\section{Functional equation, integer values, and finite products}\label{sec:functional}

The canonical product is now an independently defined object.  We derive its
recurrence rather than building the recurrence into the definition.

\subsection{Finite cancellation and the Gamma limit}

Let \(F_N(z)\) be the product in \cref{eq:Fdef} truncated after \(k=N\).
A direct cancellation gives
\begin{align}
 \frac{F_N(z)}{F_N(z-1)}
 ={}&\sqrt{2\pi}\,
 \exp\!\left(-(1+\gamma)z+\frac\gamma2-N+
                 \left(z-\frac12\right)H_N\right)\notag\\
 &\times
 \frac{(N+z)^N}{z(z+1)\cdots(z+N-1)},
 \label{eq:finite-ratio}
\end{align}
where \(H_N=\sum_{k=1}^N1/k\).  To see the only nontrivial algebra, note that
\[
 \prod_{k=1}^{N}
 \left(\frac{k+z}{k+z-1}\right)^k
 =\frac{(N+z)^N}{z\prod_{j=1}^{N-1}(z+j)}.
\]
This is a finite monomial identity and is proved by induction after clearing all
denominators.

The Gamma packet supplies the constructive Gauss limit
\begin{equation}\label{eq:gauss-gamma}
 \Gamma(z)=\lim_{N\to\infty}
 \frac{(N-1)!N^z}{z(z+1)\cdots(z+N-1)}
\end{equation}
on every compact set apart from the nonpositive integers.  Divide
\cref{eq:finite-ratio} by the finite approximation in
\cref{eq:gauss-gamma}.  The remaining factor is
\begin{equation}\label{eq:ratio-remnant}
 R_N(z)=
 \frac{\sqrt{2\pi}\,N^N}{(N-1)!N^z}
 \exp\!\left(-(1+\gamma)z+\frac\gamma2-N+
       \left(z-\frac12\right)H_N\right)
 \left(1+\frac zN\right)^N.
\end{equation}
The three standard quantitative limits are
\begin{align}
 H_N-\log N&=\gamma+O(N^{-1}),\label{eq:harmonic-limit}\\
 (N-1)!&=\sqrt{2\pi}\,N^{N-1/2}\e^{-N}(1+O(N^{-1})),
 \label{eq:stirling-fac}\\
 N\Log(1+z/N)&=z+O_R(N^{-1}).\label{eq:exp-limit}
\end{align}
Each has an explicit primitive-recursive modulus in \code{gamma\_ipra}.  On
substituting their leading terms into \cref{eq:ratio-remnant}, the powers of
\(N\) cancel, and the constant exponent is
\[
 -(1+\gamma)z+\frac\gamma2
 +(z-\tfrac12)\gamma+z=0.
\]
The quantitative error is \(O_R(1/N)\), so \(R_N(z)\to1\) uniformly on every
certified compact set.

\begin{theorem}[Functional equation]\label{thm:functional}
For every \(z\in\C\), with the identity interpreted by analytic continuation at
Gamma's poles,
\begin{equation}\label{eq:functional}
 G(z+1)=\Gamma(z)G(z),\qquad G(1)=1.
\end{equation}
On a compact zero-free domain this is a quantitative equality of fast names;
globally it is an identity of entire functions after multiplication by the
entire reciprocal-Gamma function.
\end{theorem}

\begin{proof}
Away from the nonpositive integers, combine the common product-tail modulus for
\(F_N(z)\) and \(F_N(z-1)\), the Gauss Gamma modulus, and the estimate
\(R_N(z)\to1\).  This proves \(F(z)=\Gamma(z)F(z-1)\), i.e.
\(G(z+1)=\Gamma(z)G(z)\), on a nonempty connected open domain.  Multiply by
\(1/\Gamma(z)\).  Both sides then extend to entire functions, and the equality
extends across the excluded points by their local Taylor contracts.  The
normalization is \cref{eq:G1}.
\end{proof}

For object iPRA, analytic continuation is not a black box.  At a requested
input ball intersecting a nonpositive integer, use the factorization in
\cref{thm:zeros} and the corresponding zero factor of \(1/\Gamma\).  The local
cross-multiplied identity follows from polynomial cancellation and a tail ball.
A finite cover of the input compact set combines the local proofs.

\subsection{Positive and nonpositive integer values}

Define the exact primitive-recursive sequence
\begin{equation}\label{eq:gpos}
 g_0=1,
 \qquad g_{n+1}=n!\,g_n.
\end{equation}
By \cref{eq:functional} and \(G(1)=1\),
\(G(n+1)=g_n\).  Induction yields
\begin{equation}\label{eq:integerG}
 G(n)=\prod_{j=0}^{n-2}j!
 \quad(n\ge1),
\end{equation}
where the product is empty at \(n=1\).  Thus
\[
 G(1),G(2),G(3),G(4),G(5),\ldots
 =1,1,1,2,12,288,34560,\ldots.
\]
Together with \cref{thm:zeros}, this proves
\begin{equation}\label{eq:integer-all}
 G(n)=
 \begin{cases}
 0,&n=0,-1,-2,\ldots,\\
 0!1!\cdots(n-2)!,&n=1,2,3,\ldots.
 \end{cases}
\end{equation}

\begin{theorem}[Exact iPRA integer theorem]\label{thm:integer-ipra}
Let \(\operatorname{fac}\) and \(g\) be the primitive-recursive symbols defined
by
\[
 \operatorname{fac}(0)=1,
 \quad \operatorname{fac}(n+1)=(n+1)\operatorname{fac}(n),
\]
\[
 g(0)=1,
 \quad g(n+1)=\operatorname{fac}(n)g(n).
\]
Then iPRA proves every defining equation, the superfactorial identity
\[
 g(n+1)=\prod_{j\le n}j!,
\]
and all numeral instances of the sequence above by quantifier-free induction.
\end{theorem}

\begin{proof}
The displayed recursions use the primitive-recursion constructor.  Their two
equations are instances of the zero and successor recursion axioms in
\code{PRAAx}.  Introduce a primitive-recursive accumulator for the finite
product.  Equality of its value
with \(g(n+1)\) has a quantifier-free graph formula.  The zero case is a
computation equation, and the successor case multiplies the induction
hypothesis by \((n+1)!\).  Therefore \code{PRA.qfInd} applies directly.
\end{proof}

\subsection{Superfactorials and the hyperfactorial identity}

With \(\sfac(m)=\prod_{j=0}^{m}j!\), \cref{eq:integerG} says
\begin{equation}\label{eq:superfactorial}
 G(n)=\sfac(n-2)\qquad(n\ge2),
\end{equation}
with the conventional empty extension at \(n=1\).

Let
\[
 \hyp(m)=\prod_{j=1}^{m}j^j
\]
be the hyperfactorial.  The exact finite identity on the target pages is
\begin{equation}\label{eq:hyper-cross}
 G(n)\hyp(n-1)=((n-1)!)^{n-1}\qquad(n\ge1).
\end{equation}
Indeed, count the exponent of each \(j\in\{1,\ldots,n-1\}\).  In
\(G(n)=\prod_{r=1}^{n-2}r!\), the number \(j\) appears in
\(j!, (j+1)!,\ldots,(n-2)!\), hence with exponent \(n-1-j\).  In
\(\hyp(n-1)\) it has exponent \(j\).  The total exponent is \(n-1\), which is
exactly the exponent in \(((n-1)!)^{n-1}\).

This proof can be formalized either by bounded valuation counting or, more
simply, by induction on \(n\).  The induction invariant is the division-free
identity \cref{eq:hyper-cross}; the successor step uses
\[
 G(n+1)=(n-1)!G(n),
 \quad
 \hyp(n)=\hyp(n-1)n^n,
 \quad
 n!=n(n-1)!.
\]
After substituting the induction hypothesis, both sides normalize to
\(n^n((n-1)!)^n\).  ProveIt's polynomial/power normalization modules are
designed for exactly this step.

At positive integers, \(\Gamma(n)=(n-1)!\), so
\begin{equation}\label{eq:G-over-H}
 G(n)=\frac{\Gamma(n)^{n-1}}{\hyp(n-1)}.
\end{equation}
The formal theorem is \cref{eq:hyper-cross}; the quotient form follows from the
strict positivity of \(\hyp(n-1)\).

\subsection{Finite products of Gamma values}

Iterating \cref{eq:functional} \(n\) times gives
\begin{equation}\label{eq:gamma-finite-product}
 G(k+n+1)=G(k+1)\prod_{i=1}^{n}\Gamma(k+i).
\end{equation}
Therefore, on a domain where \(G(k+1)\) is apart from zero,
\begin{equation}\label{eq:gamma-finite-quotient}
 \prod_{i=1}^{n}\Gamma(k+i)
 =\frac{G(k+n+1)}{G(k+1)}.
\end{equation}
The cross-multiplied form \cref{eq:gamma-finite-product} is global and is the
primary iPRA target.

\begin{proof}[Constructive proof]
Define
\(P_0(k)=1\) and
\(P_{n+1}(k)=P_n(k)\Gamma(k+n+1)\).  The assertion
\(G(k+n+1)=G(k+1)P_n(k)\) has the exact base case \(n=0\).  In the successor
step, multiply by \(\Gamma(k+n+1)\) and apply the functional equation once.
For named complex inputs, each multiplication is a ball operation with a
backward precision allocation.  For positive integer \(k\), every quantity is
an exact natural number, so this is ordinary quantifier-free induction.
\end{proof}

\subsection{A useful asymptotic ratio already visible from the recurrence}

For later source correction, note the exact identity for fixed \(x>0\) and
integer \(n\):
\begin{equation}\label{eq:ratio-exact}
 \frac{G(n+x)}{G(n)}
 =G(x)\,\Gamma(n)
   \prod_{j=0}^{n-1}\frac{\Gamma(x+j)}{\Gamma(j+1)}
 \qquad(n\ge1).
\end{equation}
Indeed, the extra factor \(\Gamma(n)\) exactly compensates for the last
factor introduced in the denominator product.  Gamma-ratio asymptotics, not a
power \(\Gamma(x)^n\), control the remaining product.
The precise result is proved in \cref{thm:corrected-characterization}.

\section{Taylor series, logarithmic derivatives, and power-series certificates}\label{sec:taylor}

The canonical-factor calculation \cref{eq:qseries} contains the full Taylor
series.  On \(|z|<1\), absolute convergence permits the order of summation to be
reversed:
\begin{align}
 \sum_{k=1}^{\infty}q_k(z)
 &=\sum_{k=1}^{\infty}\sum_{m=3}^{\infty}
   \frac{(-1)^{m+1}z^m}{m k^{m-1}}\notag\\
 &=\sum_{m=3}^{\infty}
   \frac{(-1)^{m+1}\zeta(m-1)}{m}z^m.
 \label{eq:q-double-sum}
\end{align}
The constructive justification is a rectangular-tail estimate.  If
\(|z|\le\rho<1\), then \(\zeta(m-1)<2\) for every \(m\ge3\), and hence all
omitted terms are dominated by a geometric series.

\begin{theorem}[Taylor series for \(\log G\)]\label{thm:taylor}
For \(|z|<1\), on the branch normalized by \(\log G(1)=0\),
\begin{equation}\label{eq:taylor-logG}
 \log G(1+z)
 =\frac{\log(2\pi)-1}{2}z
  -\frac{1+\gamma}{2}z^2
  +\sum_{m=3}^{\infty}
    \frac{(-1)^{m-1}\zeta(m-1)}{m}z^m.
\end{equation}
Equivalently, with \(k=m-1\),
\begin{equation}\label{eq:taylor-wiki}
 \log G(1+z)
 =\frac z2\log(2\pi)
  -\frac{z+(1+\gamma)z^2}{2}
  +\sum_{k=2}^{\infty}
    (-1)^k\frac{\zeta(k)}{k+1}z^{k+1}.
\end{equation}
\end{theorem}

\begin{proof}
Take the logarithm of the finite partial product in \cref{eq:Fdef} on
\(|z|<1\).  Each \(1+z/k\) lies in the disk centered at one of radius less than
one, so the Taylor branch is coherent.  Substitute \cref{eq:qseries}.  For
\(|z|\le\rho<1\),
\[
 \sum_{m=3}^{\infty}\sum_{k=1}^{\infty}
 \frac{|z|^m}{m k^{m-1}}
 \le 2\sum_{m=3}^{\infty}\frac{\rho^m}{m}<\infty.
\]
The finite rectangle identity is exact; the two tails have explicit geometric
bounds, so their difference is smaller than any requested dyadic radius.  The
linear and quadratic coefficients come from the prefactor.  The sign
\((-1)^{m-1}=(-1)^{m+1}\) is unchanged because the exponents differ by two.
\end{proof}

\subsection{Explicit truncation bounds}

Let \(T_M(z)\) contain the terms through degree \(M\), where \(M\ge3\).  If
\(|z|\le\rho<1\), then
\begin{equation}\label{eq:taylor-tail}
 \left|\log G(1+z)-T_M(z)\right|
 \le
 \frac{2\rho^{M+1}}{(M+1)(1-\rho)}.
\end{equation}
Indeed,
\[
 \sum_{m=M+1}^{\infty}
 \frac{\zeta(m-1)\rho^m}{m}
 \le \frac2{M+1}\sum_{m=M+1}^{\infty}\rho^m.
\]
For rational \(\rho=a/b<1\), the assertion that the right side is at most
\(2^{-p}\) is the natural-number inequality
\begin{equation}\label{eq:taylor-tail-cleared}
 2a^{M+1}b\,2^p
 \le (M+1)(b-a)b^{M+1}.
\end{equation}
The certificate producer increases \(M\) until this decidable inequality
holds.  Bounded search can be capped by a deliberately generous primitive-
recursive value, for example a polynomial-exponential function of
\(p+b\).  Soundness depends only on the checked inequality, not on optimality
of the cap.

The same argument gives derivative tails.  If \(0\le r\le M\), termwise
\(r\)-fold differentiation yields
\begin{equation}\label{eq:taylor-derivative-tail}
 \left|D^r\bigl(\log G(1+z)-T_M(z)\bigr)\right|
 \le
 \frac{2(M+1)^r\rho^{M+1-r}}
      {(1-\rho)^{r+1}}
\end{equation}
with a coarse but primitive-recursive constant.  A sharper falling-factorial
bound is available, but the coarse form is easier to compile and entirely
adequate.

\subsection{Exponentiating the series}

Exponentiating \cref{eq:taylor-wiki} gives
\begin{align}
 G(1+z)={}&(2\pi)^{z/2}
 \exp\!\left(-\frac{z+(1+\gamma)z^2}{2}\right)\notag\\
 &\times
 \exp\!\left(
 \sum_{k=2}^{\infty}(-1)^k\frac{\zeta(k)}{k+1}z^{k+1}
 \right).
 \label{eq:exp-taylor}
\end{align}
Comparing with the canonical product yields
\begin{equation}\label{eq:series-product-relation}
 \exp\!\left(
 \sum_{k=2}^{\infty}(-1)^k\frac{\zeta(k)}{k+1}z^{k+1}
 \right)
 =\prod_{k=1}^{\infty}E_k(z),
 \qquad |z|<1.
\end{equation}
This is not circular: both sides have independently supplied finite
approximants, and \cref{eq:q-double-sum} shows that their logarithmic
approximants differ by a rectangular tail tending effectively to zero.

For formalization, use the ball version
\[
 |u-v|\le2^{-p-c},\quad |u|,|v|\le B
 \Longrightarrow
 |\e^u-\e^v|\le \e^B2^{-p-c}.
\]
The exponential evaluator supplies an upper rational ball for \(\e^B\), and
the backward demand pass chooses \(c\).  No exact transcendental equality is
asked of the finite centers.

\subsection{The logarithmic derivative series}

Differentiate \cref{eq:taylor-logG}.  With
\[
 \phi(w):=\frac{d}{dw}\log G(w),
 \qquad
 \phi(1)=\frac{\log(2\pi)-1}{2},
\]
we obtain, for \(|w-1|<1\),
\begin{equation}\label{eq:phi-taylor}
 \phi(w)=\phi(1)-(1+\gamma)(w-1)
 +\sum_{k=2}^{\infty}(-1)^k\zeta(k)(w-1)^k.
\end{equation}
The tail on \(|w-1|\le\rho<1\) is bounded by
\begin{equation}\label{eq:phi-tail}
 \left|\sum_{k=M+1}^{\infty}(-1)^k\zeta(k)(w-1)^k\right|
 \le \frac{2\rho^{M+1}}{1-\rho}.
\end{equation}
These formulas prove the Taylor and logarithmic-derivative claims on the target
pages in a stronger complex-disk form.  One source states only \(0<z<1\);
that interval is a valid special case, not the full domain.

\subsection{Primitive-recursive names for the coefficients}

For each fixed integer \(s\ge2\), define
\[
 \zeta(s)=\sum_{n=1}^{N}\frac1{n^s}+R_{s,N},
 \qquad
 0<R_{s,N}\le\frac{1}{(s-1)N^{s-1}}.
\]
The finite sum is rational and primitive recursive.  Given \(p\), choose
\(N\) so that
\((s-1)N^{s-1}\ge2^{p+2}\).  Bounded powering and search produce a modulus.
Thus every coefficient in a fixed Taylor truncation is a fast rational name.
The development never requires a single object-language function that accepts
an arbitrary code for \(s\); the external theorem schema introduces the finite
list of fixed \(s\)'s needed at each truncation order.

Euler's constant has the standard name
\[
 \gamma_N=H_N-\log N,
 \qquad
 0<\gamma_N-\gamma<\frac1N,
\]
with a sharpened Euler--Maclaurin bound available from the Gamma packet.  The
constant \(\log(2\pi)\) is obtained from the named \(\pi\) and positive-log
evaluator.  Consequently every coefficient and every tail budget in
\cref{eq:taylor-logG} is expressible in the iPRA framework of
\cref{sec:ipra}.

\begin{roadmap}[Lean module \code{BarnesG/Taylor}]
Implement a finite coefficient function returning the rational balls for
\(\zeta(2),\ldots,\zeta(M-1)\), a Horner evaluator for \(T_M\), and the checker
for \cref{eq:taylor-tail-cleared}.  Prove the finite double-sum rearrangement by
bounded sum interchange, then use the canonical-factor tail to pass from the
finite rectangle to the product.  The object theorem is a ball inclusion, not
an equality of host-Lean real numbers.
\end{roadmap}

\section{Log-Gamma integrals, logarithmic derivatives, and the improper integral}\label{sec:integrals}

\subsection{Alexeiewsky's log-Gamma integral}

The DLMF form most convenient for branches is
\begin{equation}\label{eq:alex-dlmf}
 \Log G(1+z)=\frac z2\log(2\pi)-\frac{z(z+1)}2
 +z\Log\Gamma(1+z)-\int_0^z\Log\Gamma(1+t)\dd t.
\end{equation}
Here the logarithms are continued from the positive real axis along a path
avoiding the poles and zeros.  Replacing \(t+1\) by \(x\), or using the Gamma
functional equation, gives the equivalent target-page form
\begin{equation}\label{eq:alex}
 \int_0^z\Log\Gamma(x)\dd x
 =\frac{z(1-z)}2+\frac z2\log(2\pi)
 +(z-1)\Log\Gamma(z)-\Log G(z).
\end{equation}

\begin{theorem}[Constructive Alexeiewsky identity]\label{thm:alex}
On any rational polygonal path from a positive base point to \(z\), separated
from the nonpositive integers and equipped with coherent branches of
\(\Log\Gamma\) and \(\Log G\), \cref{eq:alex-dlmf,eq:alex} hold with an
explicit finite quadrature modulus.  Their exponentiated derivatives extend to
the corresponding meromorphic identities globally.
\end{theorem}

\begin{proof}
Define the difference between the two sides of \cref{eq:alex-dlmf} as
\(A(z)\).  Differentiate on a zero-free disk.  The canonical-product series
may be differentiated by \cref{eq:qprime-bound}.  The reciprocal-Gamma product
may be differentiated using its analogous \(k^{-2}\) tail from the Gamma
packet.  Termwise cancellation gives \(A'(z)=0\).  At \(z=0\), every term is
zero, so \(A(0)=0\).  Subdivide the certified path into disks on which both
Taylor contracts are valid.  On each segment the finite integral of the
derivative is enclosed by a rational ball of radius smaller than the allocated
budget.  A finite telescope proves \(A(z)=0\) at the requested precision.

For a more algebraic verification, start with the finite product logarithms.
After integrating the finite reciprocal-Gamma logarithm, every summand is an
explicit polynomial-log expression.  Subtract the finite Barnes logarithm;
all terms cancel except an endpoint error.  The Gamma and Barnes tail bounds
make that error tend effectively to zero.  This second route is especially
suitable for an object-language proof because it avoids a general fundamental
theorem of calculus: only a finite antiderivative identity and two tail
estimates are used.
\end{proof}

On the positive real axis, the logarithms are ordinary real logarithms and no
branch data are needed.  The integral is implemented by the fixed Gamma
quadrature contract from the existing packet.  For complex paths, the integral
is a sum over finitely many line segments and the same evaluator handles the
real and imaginary parts after an affine substitution.

\subsection{Logarithmic derivative and derivative formulas}

Differentiate \cref{eq:alex}.  The \(\log\Gamma(z)\) terms cancel and yield
\begin{equation}\label{eq:phi-global}
 \frac{G'(z)}{G(z)}
 =(z-1)\psi(z)-z+\frac12\log(2\pi)+\frac12,
\end{equation}
where \(\psi=\Gamma'/\Gamma\).  Multiplying by \(G(z)\) gives the globally
entire identity
\begin{equation}\label{eq:Gprime}
 G'(z)=G(z)\left((z-1)\psi(z)-z+\frac12\log(2\pi)+\frac12\right),
\end{equation}
understood at Gamma's poles by cancellation.  Shifting \(z\mapsto z+1\) and
using \(\psi(z+1)=\psi(z)+1/z\) gives the MathWorld form
\begin{equation}\label{eq:phi-shift}
 \frac{G'(z+1)}{G(z+1)}
 =\frac12\log(2\pi)+\frac12-z+z\frac{\Gamma'(z)}{\Gamma(z)}.
\end{equation}

\begin{proof}[Division-free formal form]
On a compact domain where \(G\) and \(\Gamma\) are apart from zero, the quotient
identities follow from the named reciprocal evaluators.  The primary theorem,
however, is obtained by multiplying \cref{eq:Gprime} by \(\Gamma(z)\):
\[
 \Gamma(z)G'(z)
 =G(z)\left((z-1)\Gamma'(z)
 +\left[-z+\tfrac12\log(2\pi)+\tfrac12\right]\Gamma(z)\right).
\]
Every factor has a Taylor-ball name, and no division appears.  Around a
nonpositive integer, use the known zero orders of \(G\) and \(1/\Gamma\) to
verify the removable singularity by finite polynomial factorization.  This is
the object-iPRA target.
\end{proof}

The Taylor series \cref{eq:phi-taylor} is recovered by differentiating the
canonical series, so the integral, product and power-series derivations are
mutually consistent without circularity.

\subsection{The improper-integral representation}

The target page gives
\begin{equation}\label{eq:improper}
 \log G(1+z)
 =\frac z2\log(2\pi)
 +\int_0^\infty\frac{1}{t}
 \left[
 \frac{1-\e^{-zt}}{4\sinh^2(t/2)}
 +\frac{z^2}{2}\e^{-t}-\frac zt
 \right]\dd t.
\end{equation}
The literal integral converges for
\begin{equation}\label{eq:improper-domain}
 \Re z>-1.
\end{equation}
At zero the three apparent singular terms cancel; at infinity the first two
terms decay exponentially and the last is \(O(t^{-2})\) after the exterior
factor \(1/t\).  If \(\Re z\le-1\), the expression may still represent the
analytic continuation, but not as this unmodified convergent improper integral.

\begin{theorem}[Improper-integral formula]\label{thm:improper}
On every compact set
\[
 |z|\le R,\qquad \Re z\ge-1+\delta,
\]
with positive rational \(\delta\), \cref{eq:improper} has a
primitive-recursive quadrature and tail modulus and equals the canonical
product definition of \(G\).
\end{theorem}

\begin{proof}
Use the exact expansion
\begin{equation}\label{eq:sinh-kernel}
 \frac{1}{4\sinh^2(t/2)}
 =\frac{\e^{-t}}{(1-\e^{-t})^2}
 =\sum_{k=1}^{\infty}k\e^{-kt},\qquad t>0.
\end{equation}
To avoid separately divergent Frullani integrals, differentiate the right side
of \cref{eq:improper} three times with respect to \(z\).  The polynomial
regularizers disappear and dominated differentiation gives
\begin{align}
 I'''(z)
 &=\int_0^\infty
   \frac{t^2\e^{-zt}}{4\sinh^2(t/2)}\dd t\notag\\
 &=\sum_{k=1}^{\infty}k\int_0^\infty t^2\e^{-(k+z)t}\dd t
 =2\sum_{k=1}^{\infty}\frac{k}{(k+z)^3}.
 \label{eq:improper-third}
\end{align}
The same third derivative follows from the product logarithm by
\cref{eq:qhigher} with \(r=3\).  The value and the first two derivatives at
\(z=0\) agree by direct finite evaluation of the regularized integral and the
coefficients in \cref{eq:taylor-logG}.  Therefore the two analytic names agree
on the connected half-plane.

For the quantitative version, split the integral into
\([0,2^{-M}]\), \([2^{-M},T]\), and \([T,\infty)\).  At zero, expand
\(1-\e^{-zt}\), \(\e^{-t}\), and
\(t^2/(4\sinh^2(t/2))\) to a fixed order.  Their constant and linear
singularities cancel exactly; the remainder is bounded by a rational multiple
of \(t\), uniformly for \(|z|\le R\).  Thus the first piece is at most
\(C_R2^{-M}\).

For the tail, \cref{eq:sinh-kernel} and
\(\Re z\ge-1+\delta\) give
\[
 \left|\frac{1-\e^{-zt}}{4\sinh^2(t/2)}\right|
 \le C_R\bigl(\e^{-t}+\e^{-\delta t}\bigr),
 \qquad t\ge1.
\]
Together with the other terms, the tail is bounded by
\[
 C_R\left(\frac{\e^{-\delta T}}{\delta T}
          +\frac{\e^{-T}}{T}+\frac1T\right).
\]
Choose primitive-recursive \(M,T\) to make the endpoint pieces smaller than
the allocated dyadic budgets.  The middle integral is ordinary rational
quadrature with a derivative modulus.  These estimates also justify the
interchanges in \cref{eq:improper-third}.
\end{proof}

\subsection{A practical evaluator}

For small \(|z|\), the Taylor series is fastest.  On moderate compact disks the
canonical product is simplest.  In the half-plane \(\Re z>-1+\delta\), the
integral can be useful when its smooth middle interval is short.  For large
\(|z|\), the asymptotic module in \cref{sec:asymptotic} dominates.  A formal
evaluator need not choose among them by an undecidable real comparison: it can
use rational upper bounds on the input box and a fixed priority table.  The
certificate records which algorithm was chosen, and the soundness theorem for
that branch supplies the result.

\begin{roadmap}[Lean modules \code{BarnesG/Integral} and
\code{BarnesG/Derivative}]
Reuse the Gamma packet's path integration and digamma contracts.  Prove the
finite version of \cref{eq:sinh-kernel} with a geometric tail.  Encode the
zero-end regularizer as a fixed Taylor polynomial plus a rational remainder.
The derivative theorem should first be cross-multiplied by \(\Gamma G\); add
quotient corollaries only under apartness structures.
\end{roadmap}

\section{Reflection, Clausen formulas, half shifts, and absolute values}\label{sec:reflection}

\subsection{Kinkelin's reflection formula}

Let
\[
 \Phi(z)=\Log G(1-z)-\Log G(1+z)+z\log(2\pi)
          -\int_0^z\pi x\cot(\pi x)\dd x.
\]
On a simply connected domain avoiding the integers and the zeros of the two
Barnes factors, differentiate and use \cref{eq:phi-global}:
\begin{align*}
 \Phi'(z)
 &=-\phi(1-z)-\phi(1+z)+\log(2\pi)-\pi z\cot(\pi z)\\
 &=z\bigl(\psi(1-z)-\psi(1+z)\bigr)+1-\pi z\cot(\pi z).
\end{align*}
The digamma recurrence and reflection formulas are
\[
 \psi(1+z)=\psi(z)+\frac1z,
 \qquad
 \psi(1-z)-\psi(z)=\pi\cot(\pi z).
\]
They give \(\Phi'(z)=0\).  Since \(\Phi(0)=0\), we obtain:

\begin{theorem}[Kinkelin reflection]\label{thm:kinkelin}
With logarithms continued coherently from \(z=0\),
\begin{equation}\label{eq:kinkelin}
 \Log G(1-z)
 =\Log G(1+z)-z\log(2\pi)
  +\int_0^z\pi x\cot(\pi x)\dd x.
\end{equation}
On a certified real interval \(-1<z<1\) this uses the ordinary real branch
away from zero crossings.  On a complex domain it is a path-dependent
logarithmic identity whose exponential is branch independent.
\end{theorem}

\begin{proof}[Constructive content]
The Gamma reflection and recurrence identities already have finite sine,
product and apartness contracts in the Gamma development.  Cover the input path
by rational disks separated from the integers.  On each disk, evaluate the
derivative of \(\Phi\) in cross-multiplied form, so no unguarded cotangent or
logarithmic derivative is used.  The resulting center is exactly zero; the
radius is made smaller than the local demand.  Integrating the zero derivative
along finitely many segments and using \(\Phi(0)=0\) gives the requested
precision.  Alternatively, expand both sides into their absolutely convergent
Taylor series on \(|z|<1\) and then continue by overlapping disks.  Every step
has a finite branch and separation certificate.
\end{proof}

\subsection{Clausen forms}

For \(0<\theta<2\pi\), define the Clausen function by either of the equivalent
constructive names
\begin{equation}\label{eq:clausen}
 \Cl(\theta)=\sum_{n=1}^{\infty}\frac{\sin(n\theta)}{n^2}
 =-\int_0^\theta\log\left(2\sin\frac t2\right)\dd t.
\end{equation}
The sine series has the uniform tail
\[
 \left|\sum_{n>N}\frac{\sin(n\theta)}{n^2}\right|\le\frac1N.
\]
On a rational compact subinterval of \((0,2\pi)\), integration of the finite
Fourier sum and the same tail proves the equality of the two names.

Integration by parts gives, for \(0<z<1\),
\begin{align}
 \int_0^z\pi x\cot(\pi x)\dd x
 &=z\log(2\sin\pi z)
   +\frac{1}{2\pi}\Cl(2\pi z).
 \label{eq:cot-clausen}
\end{align}
Indeed, the endpoint \(x\log\sin(\pi x)\) tends to zero at \(0\); this limit has
the explicit bound
\(|x\log x|\le C_p2^{-p}\) once \(x\le2^{-M(p)}\), obtained from the log
series and dyadic search.

Substitute \cref{eq:cot-clausen} into \cref{eq:kinkelin}:
\begin{equation}\label{eq:clausen-reflection}
 2\pi\log\frac{G(1-z)}{G(1+z)}
 =2\pi z\log\left(\frac{\sin\pi z}{\pi}\right)
  +\Cl(2\pi z),
 \qquad 0<z<1.
\end{equation}
Using \(G(1+z)=\Gamma(z)G(z)\), this is equivalent to
\begin{equation}\label{eq:clausen-reflection2}
 \log\frac{G(1-z)}{G(z)}
 =z\log\left(\frac{\sin\pi z}{\pi}\right)
  +\log\Gamma(z)+\frac{1}{2\pi}\Cl(2\pi z).
\end{equation}
For iPRA, the primary form is exponentiated and cross-multiplied.  The real
interval hypothesis supplies positive lower bounds for \(G(z),G(1-z)\),
\(\Gamma(z)\), and \(\sin(\pi z)\) once the input box is separated from the
endpoints.

\subsection{The half-shift tangent identity}

Put \(u=\tfrac12-z\) in \cref{eq:kinkelin}, use the Gamma reflection formula,
and simplify.  On \(|z|<\tfrac12\), where \(\cos(\pi z)>0\), one convenient
form is
\begin{equation}\label{eq:halfshift}
 \frac{G(\frac12+z)}{G(\frac12-z)}
 =\frac{(2\pi)^z}{\Gamma(\frac12+z)}
  \sqrt{\frac{\pi}{\cos(\pi z)}}
  \exp\!\left(\pi\int_0^z x\tan(\pi x)\dd x\right).
\end{equation}
An equivalent Bernoulli-polynomial form is
\begin{align}
 \log\frac{G(\frac12+z)}{G(\frac12-z)}
 ={}&\log\Gamma(\tfrac12-z)
 +(z-\tfrac12)\log(2\pi)+\tfrac12\log2\notag\\
 &+\pi\int_0^z(x-\tfrac12)\tan(\pi x)\dd x.
 \label{eq:halfshift-B1}
\end{align}
To verify equivalence, integrate
\(\frac{d}{dx}\log\cos(\pi x)=-\pi\tan(\pi x)\) and use
\[
 \Gamma(\tfrac12+z)\Gamma(\tfrac12-z)=\frac{\pi}{\cos(\pi z)}.
\]

The formal domain is not merely decorative.  A square-root evaluator needs a
positive lower bound for \(\cos(\pi z)\) on real inputs, or a branch path for
complex inputs.  On a rational interval
\(|z|\le\tfrac12-\delta\), elementary sine/cosine bounds give a primitive-
recursive positive lower bound depending on \(\delta\).  The tangent integral
is then ordinary certified quadrature.

\subsection{Conjugation and the squared absolute-value product}\label{sec:modulus}

By \cref{eq:conj},
\begin{equation}\label{eq:modsquare-basic}
 |G(z)|^2=G(z)G(\overline z).
\end{equation}
This identity is preferable to an unsquared modulus in constructive arithmetic:
it avoids square roots and sign choices.

For real \(x\notin\{0,-1,-2,\ldots\}\), divide the canonical products at
\(x+\ii y\), \(x-\ii y\), and \(x\), and simplify each conjugate pair.  The
finite algebra yields
\begin{align}
 |G(x+\ii y)|^2
 ={}&G(x)^2\exp((1+\gamma)y^2)
 \left(1+\frac{y^2}{x^2}\right)\notag\\
 &\times\prod_{k=1}^{\infty}
 \left(1+\frac{y^2}{(x+k)^2}\right)^{k+1}
 \exp\!\left(-\frac{y^2}{k}\right).
 \label{eq:modsquare}
\end{align}
The logarithm of the tail again begins at order \(k^{-2}\).  On a box
\(|x|,|y|\le R\) separated by \(\delta>0\) from the forbidden integers, expand
\(\log(1+y^2/(x+k)^2)\); after subtracting \(y^2/k\), the apparent harmonic
divergence cancels against the index shift and the finite head.  A direct
comparison with the already proved product quotient gives an explicit
\(C_{R,\delta}/N\) tail.

Taking the positive square root of \cref{eq:modsquare} gives the target-page
formula
\begin{align}
 |G(x+\ii y)|
 ={}&|G(x)|\exp\left(\frac{1+\gamma}{2}y^2\right)
 \sqrt{1+\frac{y^2}{x^2}}\notag\\
 &\times\sqrt{\prod_{k=1}^{\infty}
 \left(1+\frac{y^2}{(x+k)^2}\right)^{k+1}
 \exp\!\left(-\frac{y^2}{k}\right)}.
 \label{eq:modulus}
\end{align}
The squared identity is the primary iPRA theorem; the square-root corollary uses
nonnegativity and the positive square-root name.

At \(x=0\), extract the simple zero of \(G\) at the origin before taking the
limit.  One obtains
\begin{equation}\label{eq:pure-imag-square}
 |G(\ii y)|^2
 =y^2\exp((1+\gamma)y^2)
 \prod_{k=1}^{\infty}
 \left(1+\frac{y^2}{k^2}\right)^{k+1}
 \exp\!\left(-\frac{y^2}{k}\right),
\end{equation}
and therefore
\begin{equation}\label{eq:pure-imag}
 |G(\ii y)|
 =|y|\exp\left(\frac{1+\gamma}{2}y^2\right)
 \sqrt{\prod_{k=1}^{\infty}
 \left(1+\frac{y^2}{k^2}\right)^{k+1}
 \exp\!\left(-\frac{y^2}{k}\right)}.
\end{equation}
The absolute value on \(y\) is forced by both symmetry and nonnegativity.  This
is the corrected version of the source statement.

\subsection{Formalization pattern}

Reflection identities combine four reusable components: the derivative formula,
Gamma reflection, path integration, and a normalization point.  Clausen and
modulus formulas add only finite Fourier/product rearrangements.  A compact Lean
plan is therefore:

\begin{enumerate}
\item prove a generic ``zero derivative along a certified path'' telescope;
\item instantiate it with the cross-multiplied derivative of \(\Phi\);
\item define \(\Cl\) by its sine series and prove the integral form with a
finite Fourier sum plus the \(1/N\) tail;
\item derive \cref{eq:modsquare} directly at the level of finite partial
products, then pass to the common tail modulus;
\item expose unsquared formulas only through a positive-square-root contract.
\end{enumerate}

All induction is over finite lists, cutoffs or mesh cells and is therefore
quantifier-free after the fixed function symbols are introduced.

\section{Multiplication, duplication, and special values}\label{sec:multiplication}

\subsection{The general multiplication formula}

For an integer \(n\ge1\), put
\begin{equation}\label{eq:Pn}
 P_n(z)=\prod_{i=0}^{n-1}\prod_{j=0}^{n-1}
 G\left(z+\frac{i+j}{n}\right)
\end{equation}
and
\begin{equation}\label{eq:An}
 A_n(z)=n^{n^2z^2/2-nz}(2\pi)^{-(n^2-n)z/2}.
\end{equation}
A finite cancellation in the \(i\)-product gives
\begin{align}
 \frac{P_n(z+1/n)}{P_n(z)}
 &=\prod_{j=0}^{n-1}
   \frac{G(z+1+j/n)}{G(z+j/n)}\notag\\
 &=\prod_{j=0}^{n-1}\Gamma(z+j/n).
 \label{eq:Pshift}
\end{align}
Gauss' multiplication formula for Gamma is
\begin{equation}\label{eq:gauss-mult}
 \prod_{j=0}^{n-1}\Gamma(z+j/n)
 =(2\pi)^{(n-1)/2}n^{1/2-nz}\Gamma(nz).
\end{equation}
Also, direct subtraction of the exponents in \cref{eq:An} gives
\begin{equation}\label{eq:Ashift}
 \frac{A_n(z+1/n)}{A_n(z)}
 =n^{nz-1/2}(2\pi)^{-(n-1)/2}.
\end{equation}
Multiplying \cref{eq:Pshift,eq:gauss-mult,eq:Ashift},
\begin{equation}\label{eq:APshift}
 \frac{A_n(z+1/n)P_n(z+1/n)}{A_n(z)P_n(z)}=\Gamma(nz).
\end{equation}
The functional equation gives the same quotient for \(G(nz)\).  Hence
\begin{equation}\label{eq:periodic-ratio}
 Q_n(z):=\frac{G(nz)}{A_n(z)P_n(z)}
\end{equation}
is \(1/n\)-periodic wherever the quotient is defined; after cross
multiplication it extends through all removable points.

Apply the asymptotic expansion of \cref{sec:asymptotic} to every factor as
\(z\to+\infty\).  The coefficients of \(z^2\log z,z^2,z\log z,z,\log z\)
cancel.  The remaining constant is
\begin{equation}\label{eq:Kn}
 K(n)=\exp\bigl(-(n^2-1)\zeta'(-1)\bigr)
       n^{5/12}(2\pi)^{(n-1)/2}.
\end{equation}
Since
\begin{equation}\label{eq:A-zeta}
 \log A=\frac1{12}-\zeta'(-1),
\end{equation}
this is equivalently
\begin{equation}\label{eq:Kn-A}
 K(n)=\left(A\e^{-1/12}\right)^{n^2-1}
       n^{5/12}(2\pi)^{(n-1)/2}.
\end{equation}
Periodicity now makes the asymptotic determination exact: for fixed \(z\),
\(Q_n(z)=Q_n(z+m/n)\) for every integer \(m\), and the right side tends
effectively to \(K(n)\).

\begin{theorem}[Barnes multiplication]\label{thm:multiplication}
For every integer \(n\ge1\), with coherent branches for the powers,
\begin{equation}\label{eq:multiplication}
 G(nz)=K(n)n^{n^2z^2/2-nz}
 (2\pi)^{-(n^2-n)z/2}
 \prod_{i=0}^{n-1}\prod_{j=0}^{n-1}
 G\left(z+\frac{i+j}{n}\right).
\end{equation}
The branch-independent formal theorem is obtained by exponentiating the
quadratic logarithmic exponent along a declared path and cross multiplying all
Barnes factors.
\end{theorem}

\begin{proof}[Constructive modulus]
The finite shift calculation is exact.  Given a compact input box and precision
\(p\), choose an integer \(m\) large enough that every argument
\(z+m/n+(i+j)/n\) lies in the positive fixed sector used by the asymptotic
module and that the sum of the \(n^2+1\) explicit remainder bounds is at most
\(2^{-p-c}\).  Such an \(m\) is a primitive-recursive function of the box,
\(n\), the fixed truncation order and \(p\).  By exact periodicity, the desired
ratio at \(z\) equals the ratio at the shifted point.  The finite asymptotic
algebra reduces its ball to the constant \(K(n)\).  This proves every precision
instance.  The only use of a limit has thus been replaced by a checked finite
shift and a remainder inequality.
\end{proof}

\subsection{Duplication}

Set \(n=2\) in \cref{eq:multiplication}.  Since
\[
 P_2(x)=G(x)G(x+\tfrac12)^2G(x+1)
\]
and
\[
 K(2)=A^3\e^{-1/4}2^{5/12}(2\pi)^{1/2},
\]
solving for \(P_2(x)\) gives
\begin{equation}\label{eq:duplication}
 G(x)G(x+\tfrac12)^2G(x+1)
 =\e^{1/4}A^{-3}
 2^{-2x^2+3x-11/12}\pi^{x-1/2}G(2x).
\end{equation}
This is the Wikipedia duplication formula.  The powers of two are checked by
the exact exponent arithmetic
\[
 -2x^2+2x+x-\frac{5}{12}-\frac12
 =-2x^2+3x-\frac{11}{12}.
\]

\subsection{The value at one half}

Put \(x=1/2\) in \cref{eq:duplication}.  Since \(G(1)=1\) and
\(G(3/2)=\Gamma(1/2)G(1/2)=\sqrt\pi G(1/2)\),
\[
 \sqrt\pi\,G(1/2)^2=\e^{1/4}A^{-3}2^{1/12}.
\]
Positivity on the positive real axis selects the positive square root:
\begin{equation}\label{eq:Ghalf}
 G(1/2)=A^{-3/2}\pi^{-1/4}\e^{1/8}2^{1/24}.
\end{equation}
Two applications of the functional equation give
\begin{align}
 G(3/2)&=A^{-3/2}\pi^{1/4}\e^{1/8}2^{1/24},
 \label{eq:Gthreehalf}\\
 G(5/2)&=A^{-3/2}\pi^{3/4}\e^{1/8}2^{-23/24}.
 \label{eq:Gfivehalf}
\end{align}
These prove MathWorld equations (19)--(21).

All square roots here are positive real names.  Their iPRA proof uses lower and
upper dyadic bisection: from a positive input ball and a rational lower bound,
a bounded integer search returns the unique dyadic interval whose squared
endpoints enclose the input.

\subsection{Quarter values and Catalan's constant}

Define Catalan's constant by the alternating series
\begin{equation}\label{eq:Catalan}
 \mathcal K=\sum_{r=0}^{\infty}\frac{(-1)^r}{(2r+1)^2}.
\end{equation}
The alternating remainder is at most \((2N+3)^{-2}\), so \(\mathcal K\) has an
immediate primitive-recursive Cauchy modulus.  It is also
\(\Cl(\pi/2)\).

Let
\(a=\log G(1/4)\) and \(b=\log G(3/4)\).  Kinkelin reflection at \(z=1/4\)
gives, because \(G(5/4)=\Gamma(1/4)G(1/4)\),
\begin{equation}\label{eq:quarter-difference}
 b-a=\log\Gamma(1/4)-\frac18\log2-\frac14\log\pi
      +\frac{\mathcal K}{2\pi}.
\end{equation}
Here we used
\[
 \int_0^{1/4}\pi x\cot(\pi x)\dd x
 =\frac18\log2+\frac{\mathcal K}{2\pi}.
\]

Duplication at \(x=1/4\), followed by \cref{eq:Ghalf}, gives
\begin{equation}\label{eq:quarter-sum}
 a+b=\frac3{16}-\frac94\log A-\frac18\log2
 -\frac14\log\pi-\frac12\log\Gamma(1/4).
\end{equation}
Solving the two linear equations yields
\begin{equation}\label{eq:Gquarter}
 G(1/4)=A^{-9/8}\Gamma(1/4)^{-3/4}
 \exp\!\left(\frac3{32}-\frac{\mathcal K}{4\pi}\right),
\end{equation}
\begin{equation}\label{eq:Gthreequarter}
 G(3/4)=A^{-9/8}\Gamma(3/4)^{-1/4}
 \exp\!\left(\frac3{32}+\frac{\mathcal K}{4\pi}\right).
\end{equation}
Using
\(\Gamma(1/4)\Gamma(3/4)=\pi\sqrt2\), the second formula is equivalently
\begin{equation}\label{eq:Gthreequarter-alt}
 G(3/4)=A^{-9/8}2^{-1/8}\pi^{-1/4}
 \Gamma(1/4)^{1/4}
 \exp\!\left(\frac3{32}+\frac{\mathcal K}{4\pi}\right).
\end{equation}
Thus MathWorld's equation (18) is an alternative form of its preceding
\(G(3/4)\) value; it is not a third quarter argument.

The constructive proof does not solve a real linear system abstractly.  At
precision \(p\), it computes balls for the sum and difference with two guard
bits and applies the exact rational maps \((s,d)\mapsto((s-d)/2,(s+d)/2)\).
Exponentiation then uses a fixed positive-real evaluator.

\subsection{All odd half-integer values}

For \(k\ge2\), define
\begin{equation}\label{eq:ck}
 c_k=\prod_{i=1}^{k-2}
 \frac{2^i\Gamma(i+1/2)}{\sqrt\pi}
 =\prod_{i=1}^{k-2}(2i-1)!!.
\end{equation}
The second equality follows from
\(
\Gamma(i+1/2)=(2i)!\sqrt\pi/(4^ii!)
\), so \(c_k\) is an exact natural number.  Its primitive recursion is
\begin{equation}\label{eq:ck-rec}
 c_2=1,
 \qquad
 c_{k+1}=c_k(2k-3)!!.
\end{equation}

\begin{theorem}[Odd half-integer formula]\label{thm:odd-half}
For \(k>1\),
\begin{equation}\label{eq:odd-half}
 G(k-1/2)=
 c_k\,
 \frac{\pi^{(2k-3)/4}\e^{1/8}2^{1/24}}
      {2^{(k-1)(k-2)/2}A^{3/2}}.
\end{equation}
\end{theorem}

\begin{proof}
The case \(k=2\) is \cref{eq:Gthreehalf}.  Suppose the formula holds for
\(k\).  The functional equation gives
\[
 G(k+1/2)=\Gamma(k-1/2)G(k-1/2).
\]
Write
\[
 \Gamma(k-1/2)=\frac{(2k-3)!!}{2^{k-1}}\sqrt\pi.
\]
Multiplication by this factor changes \(c_k\) to \(c_{k+1}\), increases the
power of \(\pi\) by \(1/2\), and increases the denominator's power of two by
\(k-1\).  These are exactly the changes from \(k\) to \(k+1\) in
\cref{eq:odd-half}.  This is a quantifier-free induction once powers and
\(c_k\) are primitive-recursive symbols.
\end{proof}

\subsection{Formal dependency summary}

The multiplication theorem needs only the functional equation, Gauss
multiplication for Gamma, and the Barnes asymptotic constant.  Duplication is
its \(n=2\) instance.  The half values use duplication and recurrence.  The
quarter values use duplication, reflection and the single Clausen value
\(\Cl(\pi/2)=\mathcal K\).  Therefore the special-value module should be built
last; it contains almost no new analysis, only exact exponent arithmetic and
positive-real root/exponential operations.

\section{Hurwitz zeta, negative polygamma, and the analytic hyperfactorial}\label{sec:hurwitz}

\subsection{A constructive Hurwitz-zeta derivative}

For \(\Re z>0\), the Hurwitz zeta function has the fixed-order
Euler--Maclaurin representation
\begin{align}
 \zeta(s,z)={}&\sum_{k=0}^{N-1}(z+k)^{-s}
 +\frac{(z+N)^{1-s}}{s-1}+\frac12(z+N)^{-s}\notag\\
 &+\sum_{r=1}^{M}
 \frac{B_{2r}}{(2r)!}(s)_{2r-1}(z+N)^{-s-2r+1}
 +R_{M,N}(s,z),
 \label{eq:hurwitz-EM}
\end{align}
where \((s)_j=s(s+1)\cdots(s+j-1)\).  On a compact box in \(s,z\), with
\(\Re(z+N)\) positively bounded, the usual periodic-Bernoulli integral gives
an explicit rational majorant for \(R_{M,N}\) and its derivative in \(s\).
For fixed \(M\), all coefficients are rational; powers and logarithms use the
branch path supplied with \(z\).  Thus \(\zeta'(-1,z)\) has a
primitive-recursive fast name on every compact subset of \(\Re z>0\).

The fundamental finite shift relation is
\begin{equation}\label{eq:hurwitz-shift}
 \zeta(s,z+1)=\zeta(s,z)-z^{-s}.
\end{equation}
Differentiating at \(s=-1\) gives
\begin{equation}\label{eq:hurwitz-derivative-shift}
 \zeta'(-1,z+1)-\zeta'(-1,z)=z\Log z.
\end{equation}
The differentiation can be avoided formally: apply a symmetric rational
difference quotient in \(s\) to the exact shift identity, bound the Taylor
remainder uniformly, and then let the precision parameter increase.  This is a
schema of finite rational statements.

\subsection{Vardi's formula}

Define on \(\Re z>0\)
\begin{equation}\label{eq:vardi-rhs}
 V(z)=\frac1{12}-\zeta'(-1,z)+(z-1)\Log\Gamma(z)-\log A.
\end{equation}
Using \cref{eq:hurwitz-derivative-shift},
\(\Log\Gamma(z+1)=\Log z+\Log\Gamma(z)\), and
\(\log A=1/12-\zeta'(-1)\), one finds
\begin{align*}
 V(z+1)-V(z)
 &=-z\Log z+z\Log\Gamma(z+1)-(z-1)\Log\Gamma(z)\\
 &=\Log\Gamma(z).
\end{align*}
Also \(V(1)=0\).  The difference
\(V(z)-\Log G(z)\) is therefore one-periodic.  Euler--Maclaurin applied to
\cref{eq:vardi-rhs}, or the asymptotic theorem below, shows that this difference
tends to zero as \(\Re z\to+\infty\) in a fixed horizontal strip.  Periodicity
then forces it to be identically zero.

\begin{theorem}[Hurwitz-zeta derivative formula]\label{thm:vardi}
For \(\Re z>0\), with logarithms continued from the positive real axis,
\begin{equation}\label{eq:vardi}
 G(z)=\frac{
 \exp(1/12-\zeta'(-1,z))\,\Gamma(z)^{z-1}}{A}.
\end{equation}
Equivalently,
\[
 \Log G(z)=\frac1{12}-\zeta'(-1,z)
 +(z-1)\Log\Gamma(z)-\log A.
\]
\end{theorem}

\begin{proof}[Effective periodicity argument]
Given \(z\) in a certified compact subset of \(\Re z>0\) and precision \(p\),
choose an integer \(m\) so that the explicit Euler--Maclaurin remainders for
\(V(z+m)\) and \(\Log G(z+m)\) sum to less than \(2^{-p}\).  Repeated use of
the exact shift equations gives
\[
 V(z)-\Log G(z)=V(z+m)-\Log G(z+m).
\]
The right side lies in the dyadic ball of radius \(2^{-p}\) around zero.  This
is the precision instance of equality.  The shift count and every remainder
bound are primitive recursive in the query data.
\end{proof}

This proves MathWorld equation (4), including its domain qualification.

\subsection{Glaisher's constant as an iPRA name}

Let
\[
 \hyp(N)=\prod_{k=1}^{N}k^k.
\]
Euler--Maclaurin for \(\sum_{k=1}^N k\log k\) gives
\begin{equation}\label{eq:hyper-asymp}
 \hyp(N)=A\,
 N^{N^2/2+N/2+1/12}\e^{-N^2/4}
 \bigl(1+\varepsilon_N\bigr),
 \qquad |\varepsilon_N|\le \frac{C}{N^2}
\end{equation}
with an explicit rational \(C\) once a fixed truncation order is selected.
Thus
\begin{equation}\label{eq:A-limit}
 \log A=\lim_{N\to\infty}
 \left[
 \sum_{k=1}^{N}k\log k
 -\left(\frac{N^2}{2}+\frac N2+\frac1{12}\right)\log N
 +\frac{N^2}{4}
 \right]
\end{equation}
has a primitive-recursive modulus.  Evaluating the same Euler--Maclaurin
constant through the analytic continuation of zeta gives
\begin{equation}\label{eq:A-zeta-again}
 \log A=\frac1{12}-\zeta'(-1).
\end{equation}

For formalization one may choose \cref{eq:A-limit} as the definition of the
name \(A\), because it uses only finite sums, positive logarithms, and an
explicit remainder.  Equation \cref{eq:A-zeta-again} then becomes a theorem.
Alternatively, if the zeta module is already available, the right side of
\cref{eq:A-zeta-again} is shorter.  Both names are proved equal by comparing
their finite Euler--Maclaurin approximants.

\subsection{The analytic hyperfactorial or K-function}

To avoid confusing the K-function with the multiplication constant \(K(n)\),
write \(\mathcal H(z)\) for the analytic continuation of the hyperfactorial,
normalized by
\begin{equation}\label{eq:analytic-hyper-recurrence}
 \mathcal H(1)=1,
 \qquad
 \mathcal H(z+1)=z^z\mathcal H(z).
\end{equation}
On a branch domain define
\begin{equation}\label{eq:analytic-hyper-def}
 \mathcal H(z)=
 \frac{\exp((z-1)\Log\Gamma(z))}{G(z)}.
\end{equation}
Then \(\mathcal H(1)=1\), and
\begin{align*}
 \frac{\mathcal H(z+1)}{\mathcal H(z)}
 &=\frac{\exp(z\Log\Gamma(z+1))}{G(z+1)}
   \frac{G(z)}{\exp((z-1)\Log\Gamma(z))}\\
 &=\exp(z\Log z)=z^z.
\end{align*}
Therefore \(\mathcal H\) is the normalized analytic K-function and
\begin{equation}\label{eq:analytic-hyper-relation}
 \mathcal H(z)G(z)=\exp((z-1)\Log\Gamma(z)).
\end{equation}
At \(z=n\in\N_{>0}\), the recurrence gives
\(\mathcal H(n)=\hyp(n-1)\), and \cref{eq:analytic-hyper-relation} reduces to
the exact identity \cref{eq:hyper-cross}.

The formal primary statement is the exponentiated, cross-multiplied equality.
A global claim ``for all complex \(z\)'' must be read on the logarithmic
cover, because \(z^z\) and \(\Gamma(z)^{z-1}\) are branch dependent.  On any
fixed path domain, the continuation algorithm of \cref{sec:ipra} supplies the
needed name.

\subsection{Negative-order polygamma form}

A normalized negative-order polygamma \(\psi_{-2}\) can be characterized by
\begin{equation}\label{eq:negpoly-rec}
 \psi_{-2}(z+1)-\psi_{-2}(z)
 =z\Log z-z+\frac12\log(2\pi),
 \qquad
 \psi_{-2}(1)=\frac12\log(2\pi),
\end{equation}
plus the standard eventual higher-convexity or asymptotic normalization.  Set
\begin{equation}\label{eq:negpoly-expression}
 W(z)=-\binom z2+(z-1)\Log\Gamma(z)
      +\frac z2\log(2\pi)-\Log G(z).
\end{equation}
Using the Gamma and Barnes recurrences, a finite calculation shows that \(W\)
satisfies \cref{eq:negpoly-rec}; the asymptotic module supplies the normalization.
Hence \(W=\psi_{-2}\), and
\begin{equation}\label{eq:negative-polygamma}
 \Log G(z)=
 -\binom z2+(z-1)\Log\Gamma(z)
 +\frac z2\log(2\pi)-\psi_{-2}(z).
\end{equation}
Exponentiating gives the MathWorld form
\begin{equation}\label{eq:negative-polygamma-exp}
 G(z)=(2\pi)^{z/2}
 \exp\!\left((z-1)(\Log\Gamma(z)-z/2)-\psi_{-2}(z)\right).
\end{equation}

The qualification ``normalized'' is essential.  An indefinite second
antiderivative is determined only up to an affine function; changing that
normalization changes the right side.  The finite recurrence and value in
\cref{eq:negpoly-rec} remove the ambiguity and provide the object-iPRA target.

\begin{roadmap}[Lean modules \code{HurwitzZeta/DerivativeMinusOne} and
\code{BarnesG/Vardi}]
Implement \cref{eq:hurwitz-EM} for a fixed external Euler--Maclaurin order, with
both value and \(s\)-derivative remainder checkers.  Prove the shift relation
first at finite sums, then the derivative shift by a certified difference
quotient.  The Vardi theorem is a generic lemma: two named solutions of the
same first-order difference equation whose difference has an effective limit
zero on right shifts are equal.  Reuse that lemma for the negative-polygamma
and hyperfactorial normalizations.
\end{roadmap}

\section{Asymptotic expansion, explicit remainders, order, and type}\label{sec:asymptotic}

The symbol \(\sim\) in an asymptotic series must be replaced in iPRA by a
family of finite remainder inequalities.  For Barnes $G$, the most convenient
starting point is the Ferreira--L\'opez form with Nemes' explicit remainder.

\subsection{The Gamma-coupled expansion}

For an integer \(N\ge1\), define
\begin{align}
 \mathcal B_N(z)={}&\frac14z^2+z\Log\Gamma(z+1)
 -\left(\frac12z(z+1)+\frac1{12}\right)\Log z-\log A\notag\\
 &+\sum_{k=1}^{N-1}
 \frac{B_{2k+2}}
 {2k(2k+1)(2k+2)z^{2k}}.
 \label{eq:FL-trunc}
\end{align}
Then
\begin{equation}\label{eq:FL-expansion}
 \Log G(z+1)=\mathcal B_N(z)+R_N(z).
\end{equation}
For \(|\arg z|<\pi/2\), one exact remainder representation is
\begin{equation}\label{eq:Nemes-remainder}
 R_N(z)=\frac{(-1)^N}{2\pi^2z^{2N}}
 \int_0^\infty
 \frac{t^{2N-1}\operatorname{Li}_2(\e^{-2\pi t})}
      {1+(t/z)^2}\dd t.
\end{equation}
Equivalent periodic-Bernoulli integrals continue it to the cut plane.  They use
only a fixed Bernoulli polynomial, rational powers, and an improper integral
with exponential or algebraic tails, so they fit the framework of
\cref{sec:ipra}.

\begin{theorem}[Explicit Barnes remainder]\label{thm:Nemes-bound}
Let \(N\ge1\), \(z\ne0\), and \(\theta=\arg z\).  Then
\begin{equation}\label{eq:Nemes-bound}
 |R_N(z)|\le
 \frac{|B_{2N+2}|}
 {2N(2N+1)(2N+2)|z|^{2N}}
 \sec^{2N+1}\!\left(\frac\theta2\right),
 \qquad |\theta|<\pi.
\end{equation}
On \(|\theta|\le\pi/4\), the secant factor may be replaced by \(1\).  If
\(z>0\), \(R_N(z)\) has the sign of the first omitted term and smaller absolute
value.
\end{theorem}

\begin{proof}[Constructive reading of the remainder proof]
In \(|\theta|<\pi/2\), use \cref{eq:Nemes-remainder} and bound
\(1/|1+(t/z)^2|\) by a rational upper approximation to the appropriate angular
factor.  Expand
\(
\operatorname{Li}_2(\e^{-2\pi t})
 =\sum_{m\ge1}\e^{-2\pi mt}/m^2
\), integrate each positive term, and use the standard Bernoulli evaluation of
the resulting zeta value.  This produces exactly the coefficient in
\cref{eq:Nemes-bound}.  For wider sectors, rotate the integration path or use
the periodic-Bernoulli representation and the elementary inequality
\[
 |t+z|\ge (t+|z|)\cos(\theta/2).
\]
Raising it to the required power gives the secant factor.  On the positive
axis the integrand in a sign-normalized representation has constant sign and a
denominator strictly larger than one, proving both sign and first-omitted-term
bounds.

For a formal fixed-sector theorem, the input contains a rational
\(c>0\) with
\(
\cos(\theta/2)\ge c
\).
The verifier uses \(c^{-(2N+1)}\), avoiding an angle or inverse-trigonometric
function in the object language.  All integral cutoffs have primitive-recursive
bounds because the kernel is exponentially decaying near infinity and
polynomially regular at zero.
\end{proof}

Thus the DLMF expansion is not merely formal.  For every fixed \(N\), sector
certificate and precision, \cref{eq:Nemes-bound} gives a decidable threshold on
\(|z|\).

\subsection{The Barnes--Stirling form on the target pages}

Insert the fixed-order Stirling expansion
\begin{equation}\label{eq:stirling-gamma-use}
 \Log\Gamma(z+1)=
 \left(z+\frac12\right)\Log z-z+\frac12\log(2\pi)
 +\sum_{r=1}^{N}\frac{B_{2r}}{2r(2r-1)z^{2r-1}}
 +S_N(z)
\end{equation}
into \cref{eq:FL-expansion}.  The \(r=1\) term contributes \(1/12\).  For
\(k\ge1\), the \(r=k+1\) Gamma coefficient combines with the Barnes coefficient
as
\begin{align*}
 &\frac{B_{2k+2}}{(2k+2)(2k+1)z^{2k}}
 +\frac{B_{2k+2}}{2k(2k+1)(2k+2)z^{2k}}\\
 &\hspace{3em}=\frac{B_{2k+2}}{4k(k+1)z^{2k}}.
\end{align*}
Therefore
\begin{align}
 \Log G(z+1)={}&
 \frac{z^2}{2}\Log z-\frac{3z^2}{4}
 +\frac z2\log(2\pi)-\frac1{12}\Log z
 +\left(\frac1{12}-\log A\right)\notag\\
 &+\sum_{k=1}^{N-1}
 \frac{B_{2k+2}}{4k(k+1)z^{2k}}
 +\widetilde R_N(z),
 \label{eq:Barnes-asymptotic}
\end{align}
where
\begin{equation}\label{eq:combined-remainder}
 \widetilde R_N(z)=R_N(z)+zS_N(z).
\end{equation}
The Gamma packet supplies an explicit fixed-sector bound for \(S_N\), while
\cref{eq:Nemes-bound} bounds \(R_N\).  Hence a fully explicit bound is
\begin{align}
 |\widetilde R_N(z)|\le{}&
 \frac{|B_{2N+2}|}{|z|^{2N}}
 \left[
 \frac{C_N^{\Gamma}(\delta)}{(2N+1)(2N+2)}
 +\frac{C_N^G(\delta)}{2N(2N+1)(2N+2)}
 \right],
 \label{eq:combined-bound}
\end{align}
where \(C_N^G(\delta)\) may be taken as the rational upper bound on
\(\sec^{2N+1}(\theta/2)\), and \(C_N^\Gamma(\delta)\) is the corresponding
Stirling-remainder factor from the fixed sector
\(|\arg z|\le\pi-\delta\).  These constants are finite rational certificate
fields, not implicit big-O constants.

On the positive real axis both component remainders have the sign of their
first omitted terms and smaller magnitude.  Their first omitted coefficients
sum exactly to
\[
 \frac{|B_{2N+2}|}{4N(N+1)z^{2N}}.
\]
Thus the particularly clean bound
\begin{equation}\label{eq:positive-Barnes-bound}
 |\widetilde R_N(z)|
 <\frac{|B_{2N+2}|}{4N(N+1)z^{2N}}
 \qquad(z>0)
\end{equation}
holds, with the sign of the omitted Barnes term.  This is stronger than the
big-O statement displayed on both target pages.

Taking only the leading part of \cref{eq:Barnes-asymptotic} gives
\begin{equation}\label{eq:leading-asymptotic}
 \Log G(z+1)
 =z^2\left(\frac12\Log z-\frac34\right)
 +\frac z2\log(2\pi)-\frac1{12}\Log z+\zeta'(-1)
 +O(z^{-2}),
\end{equation}
because \(\zeta'(-1)=1/12-\log A\).  The weaker \(O(1/z)\) displayed by one
source is of course true, but the actual algebraic series contains only even
negative powers.

\subsection{An iPRA asymptotic contract}

Fix externally:
\begin{itemize}
\item a truncation order \(N\ge1\);
\item a rational sector-gap certificate \(c\le\cos(\arg z/2)\), \(c>0\);
\item a requested precision \(p\);
\item a rational lower bound \(|z|\ge L\).
\end{itemize}
The verifier computes the rational right side of
\cref{eq:Nemes-bound} or \cref{eq:combined-bound} and checks that it is at most
\(2^{-p}\) after clearing denominators.  It then checks finite balls for every
term of the truncation.  This is a quantifier-free predicate.  A producer can
increase \(L\) by bounded search until the inequality holds.  The theorem
``the asymptotic expansion is valid'' is the schema asserting soundness of this
procedure for every \(p\), not a formula containing a primitive big-O symbol.

\subsection{Order two and infinite type}

Let
\(
M_G(r)=\max_{|z|\le r}|G(z)|
\)
in classical notation.  Constructively, we need only an upper bound valid for
all named inputs in the disk and a single lower-witness input.

\begin{lemma}[Elementary global upper bound]\label{lem:order-upper}
There is an explicit integer constant \(C\) such that for every real
\(r\ge2\) and \(|z|\le r\),
\begin{equation}\label{eq:order-upper}
 \log|G(1+z)|\le C r^2\log(r+2).
\end{equation}
One may take a deliberately coarse value such as \(C=12\).
\end{lemma}

\begin{proof}
Split the canonical product at \(2r\).  For \(k\le2r\),
\[
 \log|E_k(z)|
 \le k\log(1+r/k)+r+\frac{r^2}{2k}.
\]
Sum the first term using
\(\log(1+r/k)\le\log(r+1)\) and
\(\sum_{k\le2r}k\le3r^2\).  The second contributes at most \(2r^2\), and the
third at most
\(
\frac12r^2(1+\log(2r))
\).
The logarithmic tail beyond \(2r\) is at most \(r^2/3\) by
\cref{eq:qtail}.  The prefactor is bounded by a fixed multiple of \(r^2\).
Collecting the terms gives \cref{eq:order-upper}.  Every real inequality used
has a rational majorant version.
\end{proof}

\begin{lemma}[Elementary integer lower bound]\label{lem:order-lower}
For all sufficiently large integers \(n\),
\begin{equation}\label{eq:order-lower}
 \log G(n+1)\ge \frac{n^2}{20}\log(n/2).
\end{equation}
\end{lemma}

\begin{proof}
From the exact integer formula,
\[
 \log G(n+1)=\sum_{j=1}^{n-1}(n-j)\log j.
\]
Restrict to integers \(j\) with \(n/2\le j\le3n/4\).  There are at least
\(n/5\) of them once \(n\) exceeds a fixed numeral, each satisfies
\(n-j\ge n/4\), and \(\log j\ge\log(n/2)\).  Their contribution proves
\cref{eq:order-lower}.  Floors and the small exceptional values are handled by
bounded numeral checks.
\end{proof}

\begin{theorem}[Order and type]\label{thm:order-type}
The Barnes $G$-function is entire of order exactly two and infinite type at
order two.
\end{theorem}

\begin{proof}
The upper bound implies
\[
 \limsup_{r\to\infty}\frac{\log\log M_G(r)}{\log r}\le2.
\]
The lower bound, using the witness \(z=n+1\), gives the reverse inequality.
Moreover,
\[
 \frac{\log M_G(n+1)}{(n+1)^2}
 \ge c\log(n/2)\longrightarrow\infty,
\]
so the type at order two is infinite.

The constructive replacement is explicit.  Given a positive rational
\(\varepsilon\), bounded search finds a threshold after which
\(C r^2\log(r+2)\le r^{2+\varepsilon}\); the lower witness exceeds
\(r^{2-\varepsilon}\) after another threshold.  Given an integer \(T\), choose
an explicit power of two \(n=2^{a(T+1)}\), with fixed \(a\) large enough that a
rational lower bound for \(\log2\) proves
\(\frac1{20}\log(n/2)>T\).  Then
\(\log|G(n+1)|>T(n+1)^2\) after increasing \(a\) to absorb the harmless ratio.
All thresholds are primitive recursive.
\end{proof}

\subsection{What is and is not formalized here}

The algebraic asymptotic expansion and the explicit remainder contracts are
fully developed mathematically.  The package does not claim a completed Lean
port of Nemes' periodic-Bernoulli integral.  The formalization plan is direct:
fixed Bernoulli polynomials, a rational sector lower bound, finite quadrature,
and algebraic tail integration.  Exponentially improved Stokes expansions are
outside the two target pages and are not needed for any theorem in the coverage
ledger.

\section{The two characterizations}\label{sec:characterization}

The target page gives two analogues of the Bohr--Mollerup theorem.  They require
care because a first-order arithmetic theory cannot quantify over arbitrary
real functions.  We first prove the ordinary mathematical theorems, then state
their exact coded-function interpretation.

\subsection{Third-derivative positivity for \texorpdfstring{$\log G$}{log G}}

Let \(L(x)=\log G(x)\) for \(x>0\).  From \cref{eq:phi-global},
\begin{align}
 L'''(x)&=2\psi'(x)+(x-1)\psi''(x)\notag\\
 &=2\sum_{k=0}^{\infty}
 \left(\frac1{(x+k)^2}-\frac{x-1}{(x+k)^3}\right)\notag\\
 &=2\sum_{k=0}^{\infty}\frac{k+1}{(x+k)^3}>0.
 \label{eq:Lthird}
\end{align}
For \(x\ge1\), since \(k+1\le x+k\),
\begin{equation}\label{eq:Lthird-decay}
 0<L'''(x)
 \le2\sum_{k=0}^{\infty}\frac1{(x+k)^2}
 \le2\left(\frac1{x^2}+\frac1x\right).
\end{equation}
Thus \(L'''(x)\to0\) with an explicit modulus.

The series proof is constructive.  The trigamma and second-polygamma series
have tails bounded by elementary integrals.  Positivity is witnessed by every
finite partial sum, and \cref{eq:Lthird-decay} is a rational inequality once
\(x\) is replaced by a positive rational lower bound.

\subsection{Vign\'eras uniqueness}

\begin{theorem}[Vign\'eras characterization for Barnes $G$]\label{thm:vigneras}
Let \(f:(0,\infty)\to(0,\infty)\) be three times continuously differentiable
and satisfy
\begin{equation}\label{eq:vigneras-rec}
 f(x+1)=\Gamma(x)f(x),
 \end{equation}
\begin{equation}\label{eq:vigneras-convex}
 \frac{d^3}{dx^3}\log f(x)\ge0\qquad(x\ge1).
\end{equation}
Then \(f(x)=cG(x)\) for a constant \(c>0\).  If \(f(1)=1\), then \(c=1\).
\end{theorem}

\begin{proof}
Put \(g=\log f\) and \(p=g-L\).  The two functional equations imply
\[
 p(x+1)=p(x),
\]
so \(p\) is one-periodic, as are its first three derivatives.  For every fixed
\(x\ge1\) and integer \(n\ge0\),
\[
 0\le g'''(x+n)=L'''(x+n)+p'''(x).
\]
By \cref{eq:Lthird-decay}, \(L'''(x+n)\to0\).  Therefore
\(p'''(x)\ge0\).  Since \(p'''\) is continuous and periodic,
\[
 \int_1^2p'''(x)\dd x=p''(2)-p''(1)=0.
\]
A continuous nonnegative function with zero integral is identically zero, so
\(p'''=0\).  Hence \(p\) is a polynomial of degree at most two.  A periodic
polynomial is constant.  Thus \(g=L+\log c\), and
\(f=cG\).  The normalization at one gives \(c=1\).
\end{proof}

This proof is considerably simpler than a general higher-order convexity
sandwich because the canonical solution has already been constructed and its
third derivative has an explicit decay bound.  It is also particularly well
suited to iPRA.

\begin{proof}[Finite constructive reading]
Fix a coded candidate \(f\), a rational point \(x\), and a precision \(p\).
The code supplies evaluators and continuity moduli for \(g'''\).  Choose
\(n\) so that the rational right side of \cref{eq:Lthird-decay} at \(x+n\) is
at most \(2^{-p}\).  The certified inequality \(g'''(x+n)\ge0\) and periodicity
of \(p'''\) imply
\(p'''(x)\ge-2^{-p}\).  Since this holds for every \(p\), it is the
constructive assertion \(p'''(x)\ge0\).

To infer \(p'''=0\) from zero integral, argue by a finite contradiction.  If a
ball at some rational \(x_0\) has lower endpoint \(\eta>0\), the continuity
modulus supplies a rational interval of positive length on which
\(p'''\ge\eta/2\).  A lower Riemann sum then makes
\(\int_1^2p'''>0\), contradicting the exact periodic endpoint equality.  Thus
every value ball contains zero at every requested precision.  Two certified
integrations show \(p''\), then \(p'\), are constant; periodic endpoint
equalities force those constants to be zero.  Hence \(p\) is constant.
Every step is a finite query about the fixed code.
\end{proof}

\subsection{The corrected asymptotic ratio}

The correct asymptotic law is
\begin{equation}\label{eq:correct-asymp-ratio}
 G(n+x)\sim
 \Gamma(n)^x n^{\binom{x}{2}}G(n)
 \qquad(n\to\infty),
\end{equation}
for each fixed \(x>0\).  We now prove a quantitative version.

\begin{lemma}[Digamma strip estimate]\label{lem:digamma-strip}
For every rational \(X>0\), there is an explicit rational \(C_X\) such that,
for \(n\ge2X+2\) and \(0\le t\le X\),
\begin{equation}\label{eq:digamma-strip}
 \left|
 \psi(n+t)-\log n-\frac{t-1/2}{n}
 \right|\le\frac{C_X}{n^2}.
\end{equation}
\end{lemma}

\begin{proof}
Use the first two terms of the fixed-sector digamma expansion, or differentiate
the Gamma Stirling expansion with its explicit remainder.  Alternatively,
start from
\[
 \psi(y)=\log y-\frac1{2y}
 -\int_0^\infty \frac{2u}{(u^2+y^2)(\e^{2\pi u}-1)}\dd u,
\]
whose last term is bounded by a rational multiple of \(y^{-2}\).  Expand
\(\log(n+t)=\log n+\log(1+t/n)\) to first order with the geometric remainder,
and \((n+t)^{-1}\) about \(n^{-1}\).  On \(t/n\le1/2\), all remainders have
explicit \(C_X/n^2\) bounds.
\end{proof}

Integrate the logarithmic derivative.  For \(0\le x\le X\),
\begin{equation}\label{eq:Ldifference-integral}
 L(n+x)-L(n)=\int_0^x
 \left((n+t-1)\psi(n+t)-(n+t)+c\right)\dd t,
\end{equation}
where \(c=\frac12\log(2\pi)+\frac12\).  Substitute
\cref{eq:digamma-strip} and the one-term Stirling estimate
\begin{equation}\label{eq:loggamma-one-term}
 \log\Gamma(n)=
 \left(n-\frac12\right)\log n-n+\frac12\log(2\pi)
 +O(n^{-1}).
\end{equation}
The main integrand becomes
\[
 \log\Gamma(n)+(t-\tfrac12)\log n,
\]
and the remaining error is bounded by \(D_X/n\) for an explicit rational
\(D_X\).  Integration yields
\begin{equation}\label{eq:ratio-log-error}
 \left|
 \log\frac{G(n+x)}
 {G(n)\Gamma(n)^x n^{x(x-1)/2}}
 \right|
 \le\frac{E_X}{n}.
\end{equation}
Exponentiating gives a relative-error modulus, for example
\begin{equation}\label{eq:ratio-relative-error}
 \left|
 \frac{G(n+x)}{G(n)\Gamma(n)^x n^{x(x-1)/2}}-1
 \right|
 \le \frac{2E_X}{n}
\end{equation}
whenever \(n\ge2E_X\).  This proves \cref{eq:correct-asymp-ratio} uniformly on
every fixed compact \(x\)-interval.

\begin{audit}[Why \(\Gamma(x)^n\) cannot be repaired by convention]
At \(x=1\), the correct formula is exact:
\[
 G(n+1)=\Gamma(n)G(n).
\]
Replacing \(\Gamma(n)^x\) by \(\Gamma(x)^n\) gives \(G(n+1)\sim G(n)\), whose
ratio is \(\Gamma(n)\), not one.  Therefore the source display is a genuine
transposition error, not a branch or indexing convention.
\end{audit}

\subsection{Asymptotic characterization}

\begin{theorem}[Corrected asymptotic characterization]\label{thm:corrected-characterization}
Let \(f:(0,\infty)\to(0,\infty)\) satisfy
\(f(x+1)=\Gamma(x)f(x)\).  Suppose that for every fixed \(x>0\),
\begin{equation}\label{eq:f-asymp}
 f(n+x)\sim\Gamma(n)^x n^{\binom{x}{2}}f(n).
\end{equation}
Then \(f=cG\) for a constant \(c>0\).  Conversely, every \(cG\) satisfies
\cref{eq:f-asymp}.
\end{theorem}

\begin{proof}
The converse is \cref{eq:correct-asymp-ratio}; the constant cancels.  For
uniqueness, put \(h=f/G\).  The functional equations imply
\(h(x+1)=h(x)\).  Divide the asymptotic relation for \(f\) by that for \(G\):
\[
 \frac{h(n+x)}{h(n)}\longrightarrow1.
\]
Periodicity makes the quotient independent of \(n\): it is
\(h(x)/h(1)\).  Therefore \(h(x)=h(1)\) for every \(x>0\), and \(f=cG\).
\end{proof}

The constructive version requires explicit relative-error moduli for the coded
candidate \(f\).  At requested precision \(p\), choose an integer \(n\) large
enough for the \(f\)-modulus and the explicit \(G\)-modulus
\cref{eq:ratio-relative-error}.  Positive lower bounds for the involved values
turn the two relative balls into a ball for \(h(x)/h(1)\).  Periodicity is an
exact finite iteration of the recurrence.  The resulting radius is at most
\(2^{-p}\).  Again there is no object-language quantification over all
functions; there is a theorem schema for each fixed evaluator and modulus code.

\subsection{Comparison of the two characterizations}

The derivative characterization needs a coded \(C^3\) structure and a
nonnegativity checker, but no separate asymptotic input: decay of the canonical
solution forces the periodic discrepancy to be constant.  The asymptotic
characterization needs only positive values and recurrence, but it requires a
relative-error modulus as part of the candidate data.  Both are natural in
FQNF.  The first is preferable for verifying an independently implemented
smooth function; the second is preferable when the implementation already has
an asymptotic evaluator.

\section{Lean formalization outline and path to object iPRA}\label{sec:formalization}

This section is deliberately more concrete than a list of desired theorems.  It
specifies representations, module boundaries, induction invariants and the
final theorem type.  The accompanying \code{lean/} directory mirrors this
plan.  Those files are a scaffold: the artifact-building environment did not
contain Lean, and the analytic object-theory derivations have not yet been
implemented.

\subsection{The endpoint that counts}

The target is not merely a theorem about Lean's \code{Nat}, \code{Rat}, or a
host-level complex type.  At the pinned ProveIt commit, the endpoint is
schematically
\begin{lstlisting}[language={}]
namespace LeanProofs.ConciseIPRA

-- phi : Fm PTm 0
example : [] |- [PRAAx] phi := by
  ...  -- finite object-language derivation

end LeanProofs.ConciseIPRA
\end{lstlisting}
The repository notation is \(\code{[]}\vdash_{\code{PRAAx}}\varphi\).  The syntax
\code{PR} has constructors for zero, successor, projection, composition and
primitive recursion.  \code{PRAAx} contains the defining equations, successor
axioms and quantifier-free induction, while \code{PRA.qfInd} is the convenient
induction constructor.  After an ordinary derivation is obtained, the
\code{pra\_iff\_native} bridge transfers it to the concise native scalar core.

A host theorem is useful during development but is not the final result.  Every
module should therefore have three layers:

\begin{enumerate}
\item executable host definitions and regression theorems;
\item graph formulas and bounded checkers in \code{PTm}/\code{Fm};
\item a theorem whose type is \code{Prf PRAAx [] phi}.
\end{enumerate}

\subsection{Module graph}

A practical dependency order is
\begin{verbatim}
BarnesG/FiniteProducts
  -> BarnesG/IntegerValues
  -> BarnesG/RationalBalls
  -> BarnesG/ElementaryNames
  -> BarnesG/CanonicalFactor
  -> BarnesG/CanonicalProduct
  -> BarnesG/FunctionalEquation
  -> BarnesG/Taylor
  -> BarnesG/Integral
  -> BarnesG/Derivative
  -> BarnesG/Reflection
  -> BarnesG/Multiplication
  -> HurwitzZeta/DerivativeMinusOne
  -> BarnesG/Asymptotic
  -> BarnesG/SpecialValues
  -> BarnesG/Characterization
  -> BarnesG/ObjectTargets
\end{verbatim}
The order is intentionally not the order of the source pages.  It minimizes
circularity and permits early object-level milestones.

\subsection{Stage 1: exact finite core}

Introduce primitive-recursive symbols:
\[
 \operatorname{fac},\quad g,\quad \operatorname{sf},\quad
 \operatorname{hyp},\quad \operatorname{pow},\quad
 \operatorname{facBlock},\quad c,\quad \operatorname{oddDF}.
\]
The intended equations are
\begin{align*}
 \operatorname{fac}(0)&=1,&
 \operatorname{fac}(n+1)&=(n+1)\operatorname{fac}(n),\\
 g(0)&=1,&g(n+1)&=\operatorname{fac}(n)g(n),\\
 \operatorname{hyp}(0)&=1,&
 \operatorname{hyp}(n+1)&=\operatorname{hyp}(n)(n+1)^{n+1},\\
 c(0)&=1,&c(k+1)&=c(k)\operatorname{oddDF}(k+1).
\end{align*}
All are direct \code{PR.prec} definitions.  Prove:
\begin{align}
 g(n+1)&=\operatorname{sf}(n),\label{eq:lean-sf-target}\\
 g(n)\operatorname{hyp}(n)&=\operatorname{fac}(n)^n,
 \label{eq:lean-hyp-target}\\
 g(k+n)&=g(k)\operatorname{facBlock}(k,n).
 \label{eq:lean-block-target}
\end{align}
The first and third inductions are telescopes.  The second uses the invariant
shown in \cref{eq:hyper-cross}; the successor algebra should be discharged by
the checked semiring/power normalizer in \code{IPRADerivations}, not by a large
handwritten associativity proof.

The supplied \code{lean/FiniteCore.lean} gives host definitions and the easy
telescope.  The Python checker tests 151 exact instances, but those tests are
not substitutes for \cref{eq:lean-sf-target,eq:lean-hyp-target,eq:lean-block-target}
as object derivations.

\subsection{Stage 2: rational and complex balls}

Reuse the repository's integer, rational, pairing and bounded-division
infrastructure.  A normalized fraction type is optional; an object theorem can
use raw signed numerators and positive denominators because every comparison is
cross multiplied.

Define a complex rational center as four naturals and a sign convention.  A
ball is
\[
 B(c,p)=\{z:\|z-c\|_1\le2^{-p}\}.
\]
The object predicate \code{BallEq(code,center,p)} says that the evaluator with
the fixed code returns a center whose mathematical error theorem is at most
\(2^{-p}\).  It is not necessary to encode the set braces; use a graph formula
for the evaluator and a rational inequality for the error.

Required generic lemmas are:
\begin{enumerate}
\item addition, negation, conjugation and multiplication of balls;
\item Horner evaluation with backward precision demands;
\item finite sum and product telescopes;
\item reciprocal under an apartness certificate;
\item exponential and positive logarithm Taylor evaluators;
\item positive square root by dyadic bisection;
\item path logarithm from local atanh increments.
\end{enumerate}
The Gamma packet already sketches most of these and defines host-level rational
complex arithmetic.  Each pending contract must be replaced by a bounded
checker plus a \code{Prf PRAAx} soundness theorem.

\subsection{Stage 3: named constants}

Implement fixed primitive-recursive names for:
\[
 \pi,\quad \log(2\pi),\quad \gamma,\quad
 \zeta(m)\ (m\ge2\text{ fixed}),\quad
 A,\quad \mathcal K.
\]
Recommended definitions are:
\begin{itemize}
\item Machin-type arctangent series for \(\pi\), with a geometric tail;
\item harmonic-minus-log for \(\gamma\);
\item the positive-term zeta series with integral tail;
\item the hyperfactorial Euler--Maclaurin limit \cref{eq:A-limit} for \(A\);
\item the alternating odd-square series \cref{eq:Catalan} for \(\mathcal K\).
\end{itemize}
Every name returns a rational center at precision \(p\).  Regularity is proved
by the same tail estimate that proves correctness, so no choice principle is
involved.

\subsection{Stage 4: canonical factor and product}

Define a finite rational-ball approximation to
\[
 q_k(z)=k\Log(1+z/k)-z+z^2/(2k)
\]
when \(k\ge2R\), and a direct polynomial/exponential approximation to
\(E_k(z)\) for the finite head.  The key object theorem is the cleared form of
\cref{eq:qtail}:
\begin{lstlisting}[language={}]
CanonicalTail(R,N,p) :
  2*R <= N ->
  2*R^3*2^(p+2) <= 3*N ->
  LogTailBall(R,N,p)
\end{lstlisting}
Here \code{LogTailBall} is a quantifier-free formula about the center and radius
returned by the tail evaluator.  Prove the geometric-series bound by a bounded
sum, and the \(\sum k^{-2}\) bound by the telescope
\[
 k^{-2}\le(k-1)^{-1}-k^{-1}.
\]
All divisions are cleared before induction.

Define the Barnes evaluator by
\begin{enumerate}
\item selecting \(N(R,p)\) as in \cref{eq:product-cutoff};
\item evaluating the prefactor and head factors with allocated guard bits;
\item applying the tail exponential ball;
\item multiplying the balls by a balanced tree or sequential telescope.
\end{enumerate}
A balanced tree improves numerical radii but is not logically necessary.
Soundness is an invariant over the finite factor list.

For holomorphy, implement the derivative bounds
\cref{eq:qprime-bound,eq:qsecond-bound,eq:qhigher-bound}.  It is enough initially
to formalize derivatives through order three, because all source-page theorems
can be derived from those plus fixed-order asymptotic contracts.  A later
generic order parameter can use simultaneous primitive recursion.

\subsection{Stage 5: functional equation and zeros}

The finite cancellation formula \cref{eq:finite-ratio} is a rational/power
identity after branches are exponentiated.  Reuse four Gamma contracts:

\begin{enumerate}
\item the Gauss finite product with a uniform remainder;
\item harmonic-minus-log for \(\gamma\);
\item Stirling for \((N-1)!\);
\item \((1+z/N)^N\to\e^z\) on a compact ball.
\end{enumerate}
A single telescope proves the ratio remnant \cref{eq:ratio-remnant} lies in a
ball about one.  The primary object target is
\begin{lstlisting}[language={}]
Safe(zCode,p) ->
BallEq( G(z+1), Gamma(z) * G(z), p )
\end{lstlisting}
with a local factorized version at poles.  Do not state a quotient theorem as
the primitive result.

For a fixed integer \(m\), factor the \(m\)-th canonical head polynomial as
\((z+m)^mU(z)\).  A bounded rational-box search returns a lower bound for the
other head factors, and the tail exponential is automatically apart from zero.
This produces the zero-order theorem.  Exact nonpositive-integer values are
then numeral consequences.

\subsection{Stage 6: series, integrals, and reflection}

The Taylor module uses the finite rectangle
\[
 \sum_{k\le K}\sum_{m\le M}
 \frac{(-1)^{m+1}z^m}{mk^{m-1}}
 =\sum_{m\le M}\frac{(-1)^{m+1}z^m}{m}
 \sum_{k\le K}k^{1-m}.
\]
This is a bounded-sum interchange theorem.  The two explicit tails prove
\cref{eq:taylor-tail}.  The object checker for a rational
\(\rho=a/b\) verifies \cref{eq:taylor-tail-cleared}.

For integration, reuse a generic rectangle/segment quadrature module with a
modulus field.  The log-Gamma identity is easiest to formalize from finite
products and exact antiderivatives rather than from a general theorem saying
two derivatives agree.  The improper integral needs three regions and the
regularizer described in \cref{thm:improper}.

For reflection, prove a generic finite lemma:
\begin{quote}
If a named function has a derivative ball contained in \(B(0,p+c)\) on every
cell of a certified path and its value at the initial point is zero, then its
endpoint value lies in \(B(0,p)\).
\end{quote}
Instantiate it with the cross-multiplied derivative of the Kinkelin difference.
The Clausen function is a separate sine-series name with a \(1/N\) tail.

\subsection{Stage 7: multiplication, asymptotic, and special values}

The finite product shift \cref{eq:Pshift} is a double-list cancellation.
Combine it with the already formalized Gauss multiplication theorem.  Define
the ratio \(Q_n\) only on apartness domains, but retain a cross-multiplied
periodicity theorem globally.  A right shift by a sufficiently large integer
moves every input into the asymptotic sector.  The fixed-order asymptotic
checker determines \(K(n)\), proving multiplication.

For the asymptotic module, formalize a fixed external \(N\) first.  Define the
Bernoulli polynomial \(B_{2N+2}(x)\) by a finite rational coefficient list.
Implement the periodic-Bernoulli remainder integral and prove the geometric
sector bound using a rational lower bound on \(\cos(\theta/2)\).  This suffices
for \cref{eq:Nemes-bound}.  A second theorem combines it with Gamma Stirling to
obtain \cref{eq:combined-bound}.

Special values are then short compositions.  Use a small exponent-vector
normalizer for monomials in the named positive constants
\(2,\pi,A,\Gamma(1/4),\e^{1/32},\e^{\mathcal K/(4\pi)}\).  This avoids fragile
manual rewriting of rational exponents.  The integer coefficient \(c_k\) is an
exact \code{PR} function and its half-integer theorem is qf induction.

\subsection{Stage 8: coded characterizations}

A candidate for the Vign\'eras theorem is a finite tuple of fixed programs:
\[
 (f,\log f,g',g'',g''',\text{moduli},\text{positivity bounds},
   \text{recurrence certificates}).
\]
The theorem is parameterized externally by those fixed codes.  Its premises
are object formulas saying their checkers succeed on the current finite query.
The proof in \cref{thm:vigneras} uses only right shifts, the explicit decay
bound \cref{eq:Lthird-decay}, one-period integrals, and continuity moduli.

The asymptotic characterization replaces derivative data by a relative-error
producer.  For each fixed rational \(x>0\) and precision \(p\), choose a common
\(n\) from the candidate and canonical moduli.  Periodicity then turns an
asymptotic ball into an endpoint equality ball.

\subsection{Audit and trust checks}

Each completed Lean module should finish with:
\begin{enumerate}
\item \code{\#print axioms} for host helper theorems;
\item a repository-specific linter confirming that the final theorem is a
\code{Prf PRAAx} derivation;
\item no \code{sorry}, \code{admit}, \code{axiom}, \code{Classical.choice}, or
\code{native\_decide} in the trusted path;
\item executable small-n regression examples kept logically separate from the
proof;
\item a provenance comment naming the mathematical estimate and its exact
finite version.
\end{enumerate}

The trusted endpoint is the small proof checker for the object calculus plus
Lean's kernel checking the derivation datatype.  Numerical libraries and the
Python scripts in this archive are validation aids only.

\subsection{Milestones}

A realistic sequence of independently valuable milestones is:

\begin{enumerate}
\item exact integer values, hyperfactorial relation and finite Gamma product in
\code{Prf PRAAx};
\item canonical tail and an evaluator for \(|z|\le1/2\);
\item normalization, Taylor series and the functional equation on a positive
rational domain;
\item extension to compact complex disks and zero multiplicities;
\item reflection and special half/quarter values;
\item fixed-order positive-real asymptotics;
\item fixed-sector asymptotics and multiplication;
\item the coded uniqueness theorems.
\end{enumerate}

At each stage the theorem statements remain meaningful even before the later
modules exist.  This prevents a monolithic development from hiding which
claims have actually reached object iPRA.

\section{Coverage ledger and source-by-source audit}\label{sec:coverage-audit}

This section records exactly what has been proved, in which sense, and where it
occurs.  The letters in the last column use the four status levels of
\cref{sec:scope-status}: \statusA{} is an exact arithmetic theorem already
represented by executable finite code; \statusB{} is a fully quantified
finite contract with explicit moduli; \statusC{} is a mathematically complete
constructive reduction whose routine certificate compiler remains to be
written; and \statusD{} is deliberately outside the mathematical target or is
retained only as a source annotation.  No row labelled \statusA, \statusB, or
\statusC{} is meant to assert that the final object-language derivation
\( []\vdash_{\mathrm{PRAAx}}\varphi\) has already been checked by Lean.

\renewcommand{\arraystretch}{1.17}
\small
\begin{longtable}{>{\raggedright\arraybackslash}p{0.07\textwidth}
                  >{\raggedright\arraybackslash}p{0.28\textwidth}
                  >{\raggedright\arraybackslash}p{0.40\textwidth}
                  >{\centering\arraybackslash}p{0.08\textwidth}}
\toprule
ID & Source statement & Constructive proof and qualification & Status\\
\midrule
\endfirsthead
\toprule
ID & Source statement & Constructive proof and qualification & Status\\
\midrule
\endhead
\midrule
\multicolumn{4}{r}{Continued on next page}\\
\endfoot
\bottomrule
\endlastfoot

W1 & Weierstrass canonical product for \(G(1+z)\) &
The product is defined by finite logarithmic factors.  The estimate
\(\lvert q_k(z)\rvert\le 2R^3/(3k^2)\), valid for
\(\lvert z\rvert\le R\) and \(k\ge2R\), gives a primitive-recursive cutoff,
uniform Cauchy modulus, and product enclosure; see \cref{lem:canonical-tail}.
& \statusB\\

W2 & Improper integral representation &
Direct convergence is proved for \(\Re z>-1\).  The interval is split at a
primitive-recursive truncation point; Taylor bounds near zero, exponential
tails at infinity, and rational quadrature on the middle interval give a finite
certificate.  Outside this half-plane one must use a separately certified
analytic continuation path; see \cref{thm:improper}. & \statusB\\

W3 & Entire function &
Uniform convergence of the canonical product and all fixed derivatives on
rational compact disks is reduced to explicit \(p\)-series tails.  The
constructive statement is a compact-disk family, not an appeal to a completed
space of holomorphic functions; see \cref{thm:entire}. & \statusB\\

W4 & Order two and infinite type &
The upper bound follows from the canonical product by splitting at
\(k\le2r\); the lower bound uses exact integer values and elementary factorial
bounds.  The resulting finite inequalities imply order two and infinite type;
see \cref{thm:order-type}. & \statusB\\

W5 & Functional equation \(G(z+1)=\Gamma(z)G(z)\) &
Finite products cancel exactly.  Harmonic and Stirling remainders identify the
remaining finite quotient with the Gamma Euler product.  Equality is exported
as cross-multiplied complex-ball inclusions on apartness-certified domains;
see \cref{thm:functional}. & \statusB\\

W6 & Normalization \(G(1)=1\) &
The defining product has its zero argument and every factor equal to one.  The
claim is exact and needs no limit. & \statusA\\

W7 & Integer values and superfactorials &
Quantifier-free induction on the functional equation gives
\(G(n+1)=\prod_{j<n}j!\).  The same recurrence proves the superfactorial
identification and all displayed initial values; see
\cref{thm:integer-ipra}. & \statusA\\

W8 & Hyperfactorial/K-function relation &
Both sides are reduced to finite products at integers and to the common
Gamma--Barnes recurrence for named positive-real arguments.  Division is
cleared before export to iPRA; see \cref{sec:hurwitz}. & \statusB\\

W9 & Vign\'eras third-derivative characterization &
The canonical product gives
\((\log G)'''(x)=2\sum_{k\ge0}(k+1)/(x+k)^3>0\), with an explicit tail.
For a competitor, the logarithmic quotient is periodic; positivity at large
integer translates forces its third derivative to be nonnegative, while its
period integral is zero.  Hence it is constant; see \cref{thm:vigneras}.
The coded-function compiler is still a formalization task. & \statusC\\

W10 & Duplication formula &
This is the case \(n=2\) of the multiplication theorem.  All constants are
obtained by exact finite rearrangement and the value at one half; see
\cref{sec:multiplication}. & \statusB\\

W11 & Asymptotic characterization &
The source display is corrected to
\(f(n+x)\sim\Gamma(n)^x n^{\binom{x}{2}}f(n)\).
A uniform compact-\(x\) remainder follows from the digamma strip estimate and
Euler--Maclaurin.  Periodicity of the quotient then gives uniqueness; see
\cref{thm:corrected-characterization}. & \statusB\\

W12 & Kinkelin reflection formula &
Differentiate both sides on a rational path avoiding the integers, use the
Gamma reflection identity, and fix the constant at the base point.  The
proof is formulated as path integration with a zero-free certificate; see
\cref{thm:kinkelin}. & \statusB\\

W13 & Clausen reflection forms &
Insert the finite sine-series approximation to the Clausen function in the
reflection formula.  A geometric tail bound and rational trigonometric balls
supply the modulus; see \cref{sec:reflection}. & \statusB\\

W14 & Half-shift tangent identity &
Subtract the reflection formulas at two translated arguments and use the
functional equation and elementary trigonometric identities.  The domain is a
compact interval quantitatively separated from half-integers. & \statusB\\

W15 & Taylor series for \(\log G(1+z)\) &
Termwise expansion of each finite canonical factor is justified first, then a
uniform tail permits passage to the limit.  For \(\lvert z\rvert\le\rho<1\),
the tail after degree \(M\) is at most
\(2\rho^{M+1}/((M+1)(1-\rho))\); see \cref{thm:taylor}. & \statusB\\

W16 & Exponentiated Taylor/product relation &
A rational complex exponential polynomial with its explicit tail converts the
logarithmic enclosure into an enclosure for \(G\).  No global logarithm is
asserted. & \statusB\\

W17 & General multiplication formula &
A finite double product is regrouped by residue classes modulo \(n\).  The
remaining quotient is periodic; the asymptotic formula fixes its constant.
The result is exported in exponentiated form with branch data; see
\cref{thm:multiplication}. & \statusB\\

W18 & Multiplication constant &
The constant is reduced to powers of \(2\pi\), \(n\), and the named Glaisher
constant.  For each fixed \(n\), all integer exponents and products are
primitive recursive; see \cref{sec:multiplication}. & \statusB\\

W19 & Conjugation &
Every rational finite product commutes exactly with conjugation, and the same
cutoff works for \(z\) and \(\bar z\). & \statusA\\

W20 & Absolute-value product &
Multiply the product at \(z\) by its conjugate.  The preferred statement is
for \(\lvert G(z)\rvert^2\), avoiding square roots.  Positivity and a lower
bound permit extraction of the unsquared modulus when needed. & \statusB\\

W21 & Pure-imaginary modulus &
The source prefactor \(y\) is corrected to \(\lvert y\rvert\).  The squared
identity has the exact factor \(y^2\) and is valid without a sign decision; see
\cref{audit:modulus}. & \statusB\\

W22 & Barnes asymptotic expansion &
Euler--Maclaurin with the periodic Bernoulli remainder yields the target
series.  The Nemes bound supplies a closed error term for every fixed order and
sector; see \cref{thm:Nemes-bound}. & \statusB\\

W23 & Log-Gamma integral &
Alexeiewsky's identity is proved first for finite quadrature data by
integration by parts and the functional equations, then by explicit mesh and
tail estimates; see \cref{thm:alex}. & \statusB\\

W24 & Logarithmic derivative and derivative formulas &
Differentiate the canonical logarithm on an apartness-certified disk.  The
quotient form is replaced by
\(G'=G\cdot (G'/G)\) when no lower bound for \(G\) is available. & \statusB\\

M1 & Integer hyperfactorial formula &
Finite products and the recurrence give the identity by quantifier-free
induction. & \statusA\\

M2 & Initial integer sequence &
The first values are numeral normalizations of the primitive-recursive
superfactorial.  The companion checker verifies a wider range. & \statusA\\

M3 & Hurwitz-zeta derivative formula &
Euler--Maclaurin defines a fast name for \(\zeta'(s,a)\) at \(s=-1\), with an
explicit remainder.  Subtracting the same recurrence on both sides and fixing
the value at \(a=1\) proves Vardi's formula; see \cref{thm:vardi}. & \statusB\\

M4 & Negative-order polygamma notation &
The notation is made definitional rather than imported ambiguously: the
chosen normalized antiderivative is tied to the Hurwitz-zeta derivative and a
base value.  Under this normalization the displayed Barnes formula follows
from M3. & \statusB\\

M5 & Finite Gamma product &
Repeated use of the functional equation gives a division-free product identity
with exact finite induction. & \statusA\\

M6 & Derivative formula &
This is W24 rewritten using Alexeiewsky's identity and the digamma fast name.
All divisions require an apartness witness. & \statusB\\

M7 & Leading Stirling-like asymptotic &
This is the first truncation of W22, with the remainder specialized to a closed
bound. & \statusB\\

M8 & Full Bernoulli asymptotic &
For every externally fixed truncation order, the coefficients and the Nemes
remainder bound are primitive-recursive names.  Uniformity is stated only on a
fixed rational sector. & \statusB\\

M9 & \(G(1/4)\) &
Duplication, reflection, the Gamma quarter reflection formula, and the named
Catalan and Glaisher constants determine the positive value.  Branches are
fixed by positivity. & \statusB\\

M10 & \(G(3/4)\) &
The same system determines the value; the two MathWorld forms are equivalent by
\(\Gamma(1/4)\Gamma(3/4)=\pi\sqrt2\). & \statusB\\

M11 & \(G(1/2)\) &
The duplication constant and normalization give the standard formula, with
positive-real branches. & \statusB\\

M12--M13 & \(G(3/2)\), \(G(5/2)\) &
Apply the functional equation to M11 and use the exact Gamma half-integer
values. & \statusB\\

M14 & All odd half-integer values and \(c_k\) &
Induction isolates an exact integer/rational coefficient \(c_k\) multiplying
the fixed named constants.  The recurrence for \(c_k\) is primitive recursive;
see \cref{thm:odd-half}. & \statusA/\statusB\\
\end{longtable}
\normalsize

\subsection{Items on the pages that are not additional mathematical theorems}

The historical notes, software syntax, graphs, decimal evaluations, and lists
of applications are not independent identities to prove in iPRA.  A decimal
printout is covered by a stronger statement here: a rational interval with a
checked width and containment certificate.  Likewise, an implementation note
about a computer algebra system is replaced by the explicit finite algorithms
in the archive.  The report does not attempt to formalize bibliographic claims
such as priority or naming conventions.

The asymptotic Stokes analysis beyond a fixed closed sector is also not needed
to establish the formulas on the two target pages.  It would require an
additional exponentially improved expansion and branch-dependent terminant
bounds.  Nothing in this report uses such a claim silently.

\subsection{What the ledger establishes}

The ledger covers every substantive displayed identity or characterization on
the two target pages, after the two necessary corrections.  Exact finite
identities are reduced to primitive recursion and quantifier-free induction;
analytic identities are reduced to finite rational computations with explicit
cutoffs, separation hypotheses, and error inequalities.  The remaining work is
therefore not discovery of an unspecified constructive foundation: it is the
mechanical Lean implementation of the indicated contracts and their export to
ProveIt's \code{Prf PRAAx} calculus.

\appendix
\section{Elementary estimates and certificate templates}\label{sec:appendix-estimates}

This appendix collects the inequalities used repeatedly in the main proof.  It
is deliberately elementary: each proof can be translated into bounded sums,
products, polynomial inequalities, and quantifier-free induction.  These are
the lemmas that should be formalized first because they remove nearly all
analytic ambiguity from later modules.

\subsection{A rational exponential bound}

\begin{lemma}[Finite exponential majorant]\label{lem:exp-majorant}
Let \(u\in\Q\) satisfy \(0\le u<1\).  Then
\[
  \e^u\le \frac1{1-u},
  \qquad
  0\le \e^u-1-u\le \frac{u^2}{2(1-u)}.
\]
For a complex number with \(\lvert w\rvert\le u\),
\[
  \lvert \e^w-1\rvert\le \frac{u}{1-u}.
\]
\end{lemma}

\begin{proof}
Use the power series and the factorial inequality \(j!\ge1\) for the first
bound.  For the second, \(j!\ge2\) when \(j\ge2\), hence
\[
 \sum_{j\ge2}\frac{u^j}{j!}
 \le \frac12\sum_{j\ge2}u^j
 =\frac{u^2}{2(1-u)}.
\]
The complex estimate follows from the triangle inequality.  In an iPRA
certificate the infinite geometric series is represented by the finite identity
\((1-u)\sum_{j<M}u^j=1-u^M\), together with a chosen \(M\) satisfying
\(u^M\le2^{-p}\).  Such an \(M\) is obtained by bounded search after a rational
separation witness for \(u<1\) has been supplied.
\end{proof}

\subsection{The logarithmic canonical factor}

For \(\lvert w\rvert<1\), put
\[
  r(w)=\Log(1+w)-w+\frac{w^2}{2}
      =\sum_{j=3}^{\infty}\frac{(-1)^{j-1}}{j}w^j.
\]
When \(\lvert w\rvert\le1/2\),
\begin{equation}\label{eq:r-bound-appendix}
  \lvert r(w)\rvert
  \le \sum_{j=3}^{\infty}\frac{\lvert w\rvert^j}{j}
  \le \frac{\lvert w\rvert^3}{3(1-\lvert w\rvert)}
  \le \frac{2}{3}\lvert w\rvert^3.
\end{equation}
This is the estimate used in \cref{lem:canonical-tail}.  Differentiating the
finite polynomial before taking a limit gives, for each fixed \(d\), a rational
constant \(C_d\) such that
\[
  \left|\frac{\dd^d}{\dd z^d}
   \left(k\Log(1+z/k)-z+\frac{z^2}{2k}\right)\right|
   \le \frac{C_d(1+R)^{d+3}}{k^2}.
\]
after increasing the cutoff to \(k\ge2R+2d\).  No differentiation of an
uncertified infinite object is needed.

\subsection{A telescoping bound for the quadratic tail}

For \(N\ge1\),
\begin{equation}\label{eq:pseries-tail}
  \sum_{k=N+1}^{M}\frac1{k^2}
  \le \sum_{k=N+1}^{M}\frac1{k(k-1)}
  =\frac1N-\frac1M
  \le\frac1N.
\end{equation}
The displayed equality is exact after multiplication by the common denominator
\(\prod_{k=N}^{M}k\).  Hence the estimate is especially suitable for iPRA: it
is a finite polynomial identity plus monotonicity of multiplication by positive
integers.

\subsection{Harmonic numbers and logarithms}

Let \(H_N=\sum_{k=1}^{N}1/k\).  The elementary integral comparison gives
\begin{equation}\label{eq:harmonic-log}
  0<H_N-\Log N-\gamma<\frac1N,
\end{equation}
where \(\gamma\) is introduced as the fast name
\(\gamma_N=H_N-\Log N\).  A wholly finite proof compares each rectangle with
\(\int_k^{k+1}1/t\,\dd t\); the logarithm itself is evaluated by the atanh
series after rational range reduction.  The difference
\(\gamma_N-\gamma_{2N}\) is bounded by \(1/N\), giving a dyadic Cauchy modulus.
This is the only Euler-constant input needed by the canonical product.

For fixed rational \(x\) in a compact interval and \(N\) exceeding its
apartness margin,
\begin{equation}\label{eq:log-shift}
 \Log(N+x)-\Log N
 =\sum_{j=1}^{J}\frac{(-1)^{j-1}}{j}\left(\frac{x}{N}\right)^j
  +\varepsilon_J,
 \qquad
 |\varepsilon_J|
 \le\frac{|x/N|^{J+1}}{(J+1)(1-|x/N|)}.
\end{equation}
This single formula supplies the uniform compact-\(x\) error in the corrected
asymptotic characterization.

\subsection{Products from logarithmic enclosures}

Suppose rational complex balls certify
\(|a_k|\le A_k\) and \(\sum_{k=N}^{M}A_k\le S<1\).  Then
\[
  \left|\prod_{k=N}^{M}\e^{a_k}-1\right|
  =\left|\e^{\sum a_k}-1\right|
  \le\frac{S}{1-S}.
\]
after a consistent finite branch has been chosen.  More generally, if
\(|u_k|\le1/2\), the exact finite identity
\[
 \prod_{k<N}(1+u_k)-1
 =\sum_{\varnothing\ne I\subseteq\{0,\ldots,N-1\}}
   \prod_{i\in I}u_i
\]
combined with the binomial theorem gives
\[
  \left|\prod_{k<N}(1+u_k)-1\right|
  \le \prod_{k<N}(1+|u_k|)-1.
\]
The archive's product checker uses the logarithmic version because it has a
smaller certificate, but both are primitive recursive for fixed input names.

\subsection{Safe division and logarithms}

\begin{lemma}[Ball division]\label{lem:ball-division}
Let \(z\in B(c,r)\) and \(w\in B(d,s)\), where
\(|d|-s\ge\delta>0\).  Then
\[
 \frac zw\in B\left(\frac cd,
 \frac{r}{\delta}+\frac{|c|s}{|d|\delta}\right).
\]
\end{lemma}

\begin{proof}
Write
\[
 \frac zw-\frac cd=\frac{d(z-c)+c(d-w)}{dw}
\]
and use \(|w|\ge|d|-s\ge\delta\).  Squared rational norms may replace every
square root.  Thus the checker needs only additions, multiplications,
comparisons, and one supplied lower-bound witness.
\end{proof}

A logarithm is treated similarly.  On a disk
\(B(c,r)\) with \(r<|c|\), choose a rational center logarithm and expand
\[
 \Log z=\Log c+\sum_{j\ge1}\frac{(-1)^{j-1}}j
       \left(\frac{z-c}{c}\right)^j.
\]
The tail is bounded by \cref{eq:log-shift}.  Along a path, overlapping disks
transport the branch; the finite overlap list is part of the certificate.

\subsection{Rational quadrature}

Let \(f\) have a certified derivative bound \(|f'|\le L\) on \([a,b]\).
For the left Riemann sum on a mesh of width \(h=(b-a)/m\),
\begin{equation}\label{eq:riemann-error}
 \left|\int_a^b f(t)\,\dd t-h\sum_{j<m}f(a+jh)\right|
 \le \frac{L(b-a)h}{2}.
\end{equation}
Indeed, on each interval the pointwise error is at most \(L(t-a-jh)\), and
integration gives \(Lh^2/2\).  The right side is rational when the interval and
bound are rational.  Repeated integration by parts, with all boundary terms
retained, yields higher-order rules and smaller meshes but is not logically
necessary.

For an improper integral, prove a local Taylor bound on \([0,\eta]\), use
\cref{eq:riemann-error} on \([\eta,T]\), and bound the remaining tail by a
geometric or incomplete-Gamma estimate.  Choosing \(\eta,m,T\) by bounded
search turns the whole statement into one FQNF query.

\subsection{Euler--Maclaurin with a checkable remainder}

For fixed \(m\), repeated integration by parts proves the finite identity
\begin{align}
 \sum_{k=a}^{b-1}f(k)
 &=\int_a^b f(t)\,\dd t+\frac{f(a)-f(b)}2\\
 &\quad+\sum_{j=1}^{m}
 \frac{B_{2j}}{(2j)!}
 \bigl(f^{(2j-1)}(b)-f^{(2j-1)}(a)\bigr)+R_m,
\end{align}
where
\[
 R_m=-\frac1{(2m)!}\int_a^b
 B_{2m}(\{t\})f^{(2m)}(t)\,\dd t.
\]
For externally fixed \(m\), the Bernoulli polynomial is a fixed rational
polynomial.  Its maximum on each unit interval is bounded by subdivision and
rational derivative bounds.  Thus
\[
 |R_m|\le \frac{C_m}{(2m)!}
             \int_a^b|f^{(2m)}(t)|\,\dd t
\]
is a finite checker statement.  Applied to the logarithmic summands of Barnes
\(G\), it yields the expansion of \cref{sec:asymptotic}.  Nemes' closed sector
bound is stronger and is used in the main theorem; the elementary form here is
sufficient to build the first complete Lean implementation.

\subsection{The multiplication constant by finite rearrangement}

Define the finite quotient obtained by replacing every Barnes function in the
multiplication formula by its canonical product truncated at \(M\).  Split all
indices uniquely as \(k=qn+r\), \(0\le r<n\).  The factors depending on
\(z\) cancel exactly after multiplication by the corresponding finite Gamma
multiplication identity.  What remains is independent of \(z\) up to an error
bounded by \(O_n(1/M)\), because each uncancelled logarithmic tail is bounded
by \cref{eq:pseries-tail}.  Evaluating at \(z=0\), then applying the
Euler--Maclaurin formula to the finite sums of \(k\log k\), gives the explicit
constant.  This route is preferable in iPRA to an abstract theorem saying that
a holomorphic periodic function is constant: every step is a finite residue-
class rearrangement plus a rational error estimate.

\subsection{Finite-difference form of the Vign\'eras condition}

A coded fast name need not expose a completed third derivative.  On a rational
interval and mesh \(h>0\), use
\[
 \Delta_h^3L(x)=L(x+3h)-3L(x+2h)+3L(x+h)-L(x).
\]
If \(L'''\ge0\), then \(\Delta_h^3L(x)\ge0\); conversely, a uniform family of
such inequalities with an explicit third-order Taylor remainder recovers the
third-derivative condition.  For Barnes \(G\), the product gives
\[
 \Delta_h^3L(x)
 =2h^3\sum_{k\ge0}\frac{k+1}{\xi_{k,x,h}^{3}}
\]
for suitable points in the corresponding intervals; a derivative-free proof
instead expands each logarithm and checks positivity of the resulting rational
kernel.  The tail is bounded by
\[
 2h^3\sum_{k\ge N}\frac{k+1}{(x+k)^3}
 \le 2h^3\left(\frac1{x+N-1}+\frac1{(x+N-1)^2}\right).
\]
This is the recommended object-iPRA formulation of the characterization.

\subsection{Sample object-language targets}

After fixing primitive-recursive codes for all inputs, the principal analytic
claims become ordinary first-order arithmetic sentences.  Schematically, the
canonical product target is
\[
 \forall R,p,z,N\;
 \bigl(\mathrm{Disk}(z,R)\wedge N\ge C(R,p)\bigr)
 \to
 \BallEq\bigl(P_N(z),P_{C(R,p)}(z),p\bigr),
\]
where \(z\) is a code for a rational complex point, \(P_N\) is a cleared-
denominator finite product, and \(\BallEq\) is a polynomial inequality between
squared norms.  The functional equation target is
\[
 \forall z,p\;\mathrm{Safe}(z,p)\to
 \BallEq\bigl(G_p(z+1),\Gamma_p(z)G_p(z),p\bigr).
\]
The quantifiers range only over natural numbers; every displayed predicate is a
fixed primitive-recursive symbol.  The proof in \code{Prf PRAAx} is assembled
from defining equations, order lemmas, bounded-sum/product lemmas, and
quantifier-free induction.

\subsection{Reproducible finite checks}

The companion programs do not stand in for the proofs.  They provide three
useful regression layers:
\begin{enumerate}
\item \code{verify\_finite\_core.py} checks exact factorial, superfactorial,
      hyperfactorial, finite-product, and recurrence identities on a broad
      finite range;
\item \code{certified\_product.py} uses exact rational arithmetic to compare
      canonical logarithmic tails with the proved bound;
\item \code{check\_numeric\_identities.py} independently evaluates selected
      functional, multiplication, reflection, special-value, and asymptotic
      formulas at high precision.
\end{enumerate}
Each script writes machine-readable evidence.  A failure indicates either a
formula transcription error or an implementation bug; success is only a
regression check, never a substitute for universal induction or an analytic
remainder theorem.

\subsection{Minimal trusted base}

The intended final trust chain is:
\[
 \begin{gathered}
 \text{Lean kernel}
 \longleftarrow \code{Prf PRAAx}\text{ derivations}
 \longleftarrow \text{certificate checkers},\\[2mm]
 \text{certificate checkers}
 \longleftarrow \text{explicit rational certificates}.
 \end{gathered}
\]
The host Python programs and the prose report are not in the trusted base.
Neither is Mathlib analysis: it may be used to help construct proofs, but the
endpoint is the object calculus already present in ProveIt.  This separation is
what makes the phrase ``proved in iPRA'' precise rather than rhetorical.


\begin{thebibliography}{99}
\bibitem{Barnes1900}
E. W. Barnes, \emph{The theory of the $G$-function}, Quarterly Journal of Pure and Applied Mathematics 31 (1900), 264--314.

\bibitem{DLMF}
NIST Digital Library of Mathematical Functions, \emph{Section 5.17: Barnes' $G$-Function (Double Gamma Function)},
\url{https://dlmf.nist.gov/5.17}, inspected 16 November 2026.

\bibitem{Wikipedia}
Wikipedia contributors, \emph{Barnes $G$-function},
\url{https://en.wikipedia.org/wiki/Barnes_G-function}, inspected 16 November 2026.

\bibitem{MathWorld}
E. W. Weisstein, \emph{Barnes $G$-Function}, MathWorld---A Wolfram Web Resource,
\url{https://mathworld.wolfram.com/BarnesG-Function.html}, inspected 16 November 2026.

\bibitem{Nemes2014}
G. Nemes, \emph{Error bounds and exponential improvement for the asymptotic expansion of the Barnes $G$-function},
\url{https://arxiv.org/abs/1406.2535}, 2014.

\bibitem{MarichalZenaidi}
J.-L. Marichal and N. Zena\"idi, \emph{A Generalization of Bohr--Mollerup's Theorem for Higher Order Convex Functions}, Springer, 2022; especially the chapter ``Applications to Some Standard Special Functions.''

\bibitem{Adamchik}
V. S. Adamchik, papers on the Barnes function, multiple gamma functions and related constants, as cited by MathWorld and DLMF.

\bibitem{ProveIt}
V. Reshetnikov, \emph{ProveIt}, commit
\code{2a7a1cc230762eaa57df921e944a5443afeb0519},
\url{https://github.com/VladimirReshetnikov/ProveIt}.

\bibitem{GammaPacket}
ProveIt, \emph{Gamma in intuitionistic primitive recursive arithmetic},
\code{Logic/PrimitiveRecursiveArithmetic/docs/proofs/analysis/gamma\_ipra},
commit as in \cite{ProveIt}.
\end{thebibliography}

\end{document}
